% !TeX encoding = UTF-8
% !TeX spellcheck = en_US
% !TeX root = lab-manual.tex
% !TeX TXS-program:compile = txs:///xelatex/[--shell-escape]

\documentclass{book}

\usepackage{tabularx}

\usepackage[utf8]{inputenc}
\usepackage{amsmath}
\usepackage{amssymb}
\usepackage{listings}
\usepackage{xcolor}
\usepackage{booktabs}
\usepackage{geometry}
\usepackage{hyperref}

\geometry{a4paper, margin=1in}

\definecolor{codegreen}{rgb}{0,0.6,0}
\definecolor{codegray}{rgb}{0.5,0.5,0.5}
\definecolor{codepurple}{rgb}{0.58,0,0.82}
\definecolor{backcolour}{rgb}{0.95,0.95,0.92}

\lstdefinestyle{mystyle}{
        backgroundcolor=\color{backcolour},   
        commentstyle=\color{codegreen},
        keywordstyle=\color{magenta},
        numberstyle=\tiny\color{codegray},
        stringstyle=\color{codepurple},
        basicstyle=\ttfamily\small,
        breakatwhitespace=false,         
        breaklines=true,                 
        captionpos=b,                    
        keepspaces=true,                 
        numbers=left,                    
        numbersep=5pt,                  
        showspaces=false,                
        showstringspaces=false,
        showtabs=false,                  
        tabsize=2
}

\lstset{style=mystyle}


\usepackage{blindtext}
\usepackage{titlesec}
\title{lab.bitsmasher.net Manual}
\author{franklin}
\date{August 2026 }

\begin{document}

\maketitle

\tableofcontents

\chapter{Introduction}

\section{Purpose and Scope}
This manual documents the lab.bitsmasher.net infrastructure -- network, hardware, services, and procedures. It is intended as a reference for current operations and a foundation for onboarding new administrators.

% \input{sections/history}
% (historical content kept as legacy; see individual sections)

\chapter{Hardware Inventory}

\section{Desk Setup}
Franklin's desk hosts chonk, the primary gateway workstation. Sly's desk hosts a GAMING Windows PC at the main office near music.lan's monitor. Both have access to all lab infrastructure.

% \section{Hardware Inventory}

\subsection{Desk Setup}

\begin{itemize}
    \item \textbf{Franklin's desk}: chonk (10.10.8.60) -- GAMING Windows PC, primary workstation, Gateway host running OpenClaw
    \item \textbf{Sly's desk}: GAMING Windows computer (Twitch: \url{https://www.twitch.tv/s1y_b0rg}) -- at main office near music.lan monitor
    \item Both sit at the main office where music.lan's monitor is displayed
\end{itemize}

\subsection{Server and Compute Hosts}

\subsubsection*{Core Infrastructure (Debian 12)}

\begin{tabularx}{\textwidth}{l l l X}
\textbf{Host} & \textbf{IP} & \textbf{User} & \textbf{Role} \\
\hline
chonk & 10.10.8.60 & franklin/openclaw & Gateway host, primary workstation\\
stargate & 10.10.16.66 & openclaw & Ansible collection workspace, build host\\
skynet & 10.10.16.10 & openclaw/franklin & Minecraft server, IP-direct login (hostname resolves oddly)\\
time & 10.10.12.2 & openclaw/franklin & Stratum-1 NTP server with GPS reference\\
wonderland & 178.62.60.55 & franklin/openclaw & Public cloud VM (Debian 12, 6.1 kernel)\\
music.lan & 192.168.86.38 & root & Desktop machine via skynet route
\end{tabularx}

\subsubsection*{Jetson Cluster (\textasciitilde 87 days uptime)}

All three Jetsons run Ubuntu 18.04 with ed25519 key auth (openclaw permanent, franklin has password \texttt{123} as fallback):

\begin{tabularx}{\textwidth}{l l X}
\textbf{Host} & \textbf{IP} & \textbf{Status} \\
\hline
node900 & 10.10.12.90 & Online, key auth\\
node901 & 10.10.12.91 & Online, key auth\\
node903 & 10.10.12.93 & Online, key auth\\
node902 & 10.10.12.92 & Off / no route to host (likely powered down)
\end{tabularx}

\subsubsection*{KDC and Directory Services}

\begin{tabularx}{\textwidth}{l l X}
\textbf{Host} & \textbf{IP} & \textbf{Status} \\
\hline
odroid-c1 / kdc1 & 10.10.12.254 & Online (~242d uptime), root key auth, KDC for lab.bitsmasher.net Kerberos realm\\
ldap/bbb1 & 10.10.13.1 & SSH reachable as root; slapd failed Dec 2025 (TLS cert issue)
\end{tabularx}

\subsection{Historical / Decommissioned Hosts}

\begin{itemize}
    \item \textbf{snowy}: Hard disk physically removed. /mnt/snowy is now just a mounted drive (not SSH-accessible). Should be cleaned from DNS and SSH config.
    
    \item \textbf{blowfish}: Previously at 10.10.12.15, reassigned to 10.10.14.85. Key not authorized yet -- unreachable from chonk but port 22 is open. Expected role: database host with OpenBSD.
    
    \item \textbf{node902}: Powered off; no route to host from any known peer.
\end{itemize}

\subsection{Planned Hardware Expansion}

\begin{itemize}
    \item blowfish as the primary database server (OpenBSD planned)
    \item head2 for cluster management (referenced in network\_update.sh but unreachable)
    \item k3s_server and k3s_agent roles for Kubernetes across node infrastructure
\end{itemize}

\section{Hardware Inventory}

\subsection{Desk Setup}

\begin{itemize}
    \item \textbf{Franklin's desk}: chonk (10.10.8.60) -- GAMING Windows PC, primary workstation, Gateway host running OpenClaw
    \item \textbf{Sly's desk}: GAMING Windows computer (Twitch: \url{https://www.twitch.tv/s1y_b0rg}) -- at main office near music.lan monitor
    \item Both sit at the main office where music.lan's monitor is displayed
\end{itemize}

\subsection{Server and Compute Hosts}

\subsubsection*{Core Infrastructure (Debian 12)}

\begin{tabularx}{\textwidth}{l l l X}
\textbf{Host} & \textbf{IP} & \textbf{User} & \textbf{Role} \\
\hline
chonk & 10.10.8.60 & franklin/openclaw & Gateway host, primary workstation\\
stargate & 10.10.16.66 & openclaw & Ansible collection workspace, build host\\
skynet & 10.10.16.10 & openclaw/franklin & Minecraft server, IP-direct login (hostname resolves oddly)\\
time & 10.10.12.2 & openclaw/franklin & Stratum-1 NTP server with GPS reference\\
wonderland & 178.62.60.55 & franklin/openclaw & Public cloud VM (Debian 12, 6.1 kernel)\\
music.lan & 192.168.86.38 & root & Desktop machine via skynet route
\end{tabularx}

\subsubsection*{Jetson Cluster (\textasciitilde 87 days uptime)}

All three Jetsons run Ubuntu 18.04 with ed25519 key auth (openclaw permanent, franklin has password \texttt{123} as fallback):

\begin{tabularx}{\textwidth}{l l X}
\textbf{Host} & \textbf{IP} & \textbf{Status} \\
\hline
node900 & 10.10.12.90 & Online, key auth\\
node901 & 10.10.12.91 & Online, key auth\\
node903 & 10.10.12.93 & Online, key auth\\
node902 & 10.10.12.92 & Off / no route to host (likely powered down)
\end{tabularx}

\subsubsection*{KDC and Directory Services}

\begin{tabularx}{\textwidth}{l l X}
\textbf{Host} & \textbf{IP} & \textbf{Status} \\
\hline
odroid-c1 / kdc1 & 10.10.12.254 & Online (~242d uptime), root key auth, KDC for lab.bitsmasher.net Kerberos realm\\
ldap/bbb1 & 10.10.13.1 & SSH reachable as root; slapd failed Dec 2025 (TLS cert issue)
\end{tabularx}

\subsection{Historical / Decommissioned Hosts}

\begin{itemize}
    \item \textbf{snowy}: Hard disk physically removed. /mnt/snowy is now just a mounted drive (not SSH-accessible). Should be cleaned from DNS and SSH config.
    
    \item \textbf{blowfish}: Previously at 10.10.12.15, reassigned to 10.10.14.85. Key not authorized yet -- unreachable from chonk but port 22 is open. Expected role: database host with OpenBSD.
    
    \item \textbf{node902}: Powered off; no route to host from any known peer.
\end{itemize}

\subsection{Planned Hardware Expansion}

\begin{itemize}
    \item blowfish as the primary database server (OpenBSD planned)
    \item head2 for cluster management (referenced in network\_update.sh but unreachable)
    \item k3s_server and k3s_agent roles for Kubernetes across node infrastructure
\end{itemize}


\chapter{Network Infrastructure}

\section{Addressing and Subnets}
The lab spans multiple private subnets organized by function, from the main office (10.10.8.0/21) to the backbone (10.10.16.0/) and specialized service zones.

% \section{Network Infrastructure}

\subsection{Addressing Scheme}

The lab uses private IPv4 addressing across multiple subnets, organized by function:

\begin{tabularx}{\textwidth}{l l X}
\textbf{Subnet} & \textbf{Allocation} & \textbf{Purpose} \\
\hline
10.10.8.0/21 & chonk, chonk-wifi, music.lan (via skynet) & Main office / desk area\\
10.10.12.0/& time, node90x, odroid-c1, blowfish, ns1 & Core infrastructure\\
10.10.13.0/ & ldap/bbb1 & Directory services (moved to working hosts)\\
10.10.14.0/ & blowfish (reassigned) & New allocation for blowfish\\
10.10.15.0/ & femputer & Guest / temporary hosts\\
10.10.16.0/& stargate, skynet, edge-t & Lab backbone / staging
\end{tabularx}

Note: The full subnet masks vary by role and DHCP scope configuration on the Dream Machine (DHCP server).

\subsection{DHCP}

The \texttt{dhcp} Ansible role manages the DHCP server configuration stored at \texttt{roles/dhcp/files/dhcpd.conf}. Key reserved hosts include:

\begin{itemize}
    \item \textbf{chonk}: 10.10.8.60, MAC fc:9d:05:01:27:02 -- Franklin's desk
    \item \textbf{chonk-wifi}: 10.10.8.61 -- wireless client
    \item \textbf{chonk-ten-gig-1}: 10.10.8.?? -- dedicated ten-gig link
    \item \textbf{blowfish}: 10.10.8.15 / 10.10.12.15, MAC 98:b7:85:21:ad:77 -- database host
    \item \textbf{femputer}: 10.10.15.1, MAC d8:bb:c1:af:21:17 -- temporary/guest
\end{itemize}

\subsection{DNS}

The DNS infrastructure is managed by the \texttt{dns} Ansible role (BIND/named). Zones maintained include:
\begin{itemize}
    \item Forward zone: db.home.lab
    \item Reverse zones for each subnet
    \item A records and PTR records for all lab hosts
\end{itemize}

DNS role molecule testing confirms BIND packages install correctly and named.conf.options are written. The DNS server (ns1) has been offline during recent audits, which explains why many hosts in the 10.10.x.0/24 range cannot resolve -- ns1 is a single point of failure for lab DNS resolution.

\subsection{Firewall and Access Control}

\begin{itemize}
    \item The Dream Machine provides built-in DHCP and handles basic routing between subnets
    \item odroid-c1 serves as the KDC and likely has pf firewall rules managing cross-subnet Kerberos traffic
    \item OpenBSD hosts (blowfish planned) use pf for host-level firewalls
    \item stargate has configured SSH access from chonk via openclaw key pair
\end{itemize}

\subsection{Physical Topology}

The lab spans multiple physical locations and subnets:
\begin{itemize}
    \item Main office where Franklin and Sly have their desks monitors (music.lan)
    \item Server racks hosting core infrastructure (time, odroid-c1/KDC)
    \item Jetson cluster (node900--node903) for GPU compute
    \item Public cloud (wonderland/178.62.60.55) for external-facing services
\end{itemize}

Routing between the 192.168.86.x subnet (music) and the rest of the lab requires hopping through skynet due to the non-standard routing configuration.

\section{Network Infrastructure}

\subsection{Addressing Scheme}

The lab uses private IPv4 addressing across multiple subnets, organized by function:

\begin{tabularx}{\textwidth}{l l X}
\textbf{Subnet} & \textbf{Allocation} & \textbf{Purpose} \\
\hline
10.10.8.0/21 & chonk, chonk-wifi, music.lan (via skynet) & Main office / desk area\\
10.10.12.0/& time, node90x, odroid-c1, blowfish, ns1 & Core infrastructure\\
10.10.13.0/ & ldap/bbb1 & Directory services (moved to working hosts)\\
10.10.14.0/ & blowfish (reassigned) & New allocation for blowfish\\
10.10.15.0/ & femputer & Guest / temporary hosts\\
10.10.16.0/& stargate, skynet, edge-t & Lab backbone / staging
\end{tabularx}

Note: The full subnet masks vary by role and DHCP scope configuration on the Dream Machine (DHCP server).

\subsection{DHCP}

The \texttt{dhcp} Ansible role manages the DHCP server configuration stored at \texttt{roles/dhcp/files/dhcpd.conf}. Key reserved hosts include:

\begin{itemize}
    \item \textbf{chonk}: 10.10.8.60, MAC fc:9d:05:01:27:02 -- Franklin's desk
    \item \textbf{chonk-wifi}: 10.10.8.61 -- wireless client
    \item \textbf{chonk-ten-gig-1}: 10.10.8.?? -- dedicated ten-gig link
    \item \textbf{blowfish}: 10.10.8.15 / 10.10.12.15, MAC 98:b7:85:21:ad:77 -- database host
    \item \textbf{femputer}: 10.10.15.1, MAC d8:bb:c1:af:21:17 -- temporary/guest
\end{itemize}

\subsection{DNS}

The DNS infrastructure is managed by the \texttt{dns} Ansible role (BIND/named). Zones maintained include:
\begin{itemize}
    \item Forward zone: db.home.lab
    \item Reverse zones for each subnet
    \item A records and PTR records for all lab hosts
\end{itemize}

DNS role molecule testing confirms BIND packages install correctly and named.conf.options are written. The DNS server (ns1) has been offline during recent audits, which explains why many hosts in the 10.10.x.0/24 range cannot resolve -- ns1 is a single point of failure for lab DNS resolution.

\subsection{Firewall and Access Control}

\begin{itemize}
    \item The Dream Machine provides built-in DHCP and handles basic routing between subnets
    \item odroid-c1 serves as the KDC and likely has pf firewall rules managing cross-subnet Kerberos traffic
    \item OpenBSD hosts (blowfish planned) use pf for host-level firewalls
    \item stargate has configured SSH access from chonk via openclaw key pair
\end{itemize}

\subsection{Physical Topology}

The lab spans multiple physical locations and subnets:
\begin{itemize}
    \item Main office where Franklin and Sly have their desks monitors (music.lan)
    \item Server racks hosting core infrastructure (time, odroid-c1/KDC)
    \item Jetson cluster (node900--node903) for GPU compute
    \item Public cloud (wonderland/178.62.60.55) for external-facing services
\end{itemize}

Routing between the 192.168.86.x subnet (music) and the rest of the lab requires hopping through skynet due to the non-standard routing configuration.


\section{DNS}
BIND on ns1 provides authoritative DNS for the lab, with forward and reverse zones plus Kerberos SRV records. Currently offline -- see Disaster Recovery for recovery procedures.

% \section{DNS Configuration}

\subsection{BIND Overview}

The lab uses BIND (named) as its authoritative DNS server. The DNS infrastructure is managed by the \texttt{dns} Ansible role in the lab-franklin collection, located at:
\begin{lstlisting}[style=mystyle]
ansible/collections/ansible_collections/lab/franklin/roles/dns/
    \end{lstlisting}

\subsection{Zones Maintained}

\begin{itemize}
    \item \textbf{Forward zone}: db.home.lab -- maps hostnames to IPs for all lab infrastructure
    
    \item \textbf{Reverse zones}: PTR records for each subnet (10.10.x.0/24 ranges)
    
    \item \textbf{SRV records}: Kerberos service discovery (\_kerberos.\_tcp, \_kdc.tcp)
\end{itemize}

\subsection{DNS Server Location}

The DNS server (ns1) is the single authoritative nameserver for the lab. It has been unreachable during recent infrastructure audits, which explains why many hosts cannot resolve hostnames -- ns1 \textit{is} the DNS server and it's offline.

When ns1 is available, it resolves all internal hostnames. Key entries include:
\begin{itemize}
    \item time.lab.bitsmasher.net $\rightarrow$ 10.10.12.2 (NTP server)
    \item odroid-c1.lab.bitsmasher.net $\rightarrow$ 10.10.12.254 (KDC)
    \item ldap.lab.bitsmasher.net $\rightarrow$ 10.10.13.1 (OpenLDAP)
    \item blowfish.lab.bitsmasher.net $\rightarrow$ 10.10.14.85 (reassigned from 10.10.12.15)
\end{itemize}

\subsection{Molecule Test Harness}

The DNS role has a molecule test harness that confirms:
\begin{enumerate}
    \item \textbf{Converge}: BIND packages installed, /etc/bind directory created, named.conf.options written
    \item \textbf{Test sequence}: dependency $\rightarrow$ syntax $\rightarrow$ create (podman container) $\rightarrow$ prepare $\rightarrow$ converge
    
    \item \textbf{Limitation}: molecule test uses Ubuntu 20.04 container with hostname "server1"; the role has a hostname gate that targets the real DNS server's hostname
\end{enumerate}

The DNS molecule test confirms correct package installation and configuration file placement but does not verify actual name resolution (no live DNS server in the container test environment).

\subsection{DNS Failure Impact}

When ns1 is offline:
\begin{itemize}
    \item Lab hosts fall back to /etc/hosts entries (which are authoritative but not dynamic)
    
    \item Kerberos KDC discovery via SRV records fails -- clients need manual kdc = entry in krb5.conf
    
    \item DHCP reservations on the Dream Machine remain functional (DHCP server is separate from DNS)
\end{itemize}

The DNS role's molecule test harness ensures that when the server comes back online, its configuration will install and write correctly.

\section{DNS Configuration}

\subsection{BIND Overview}

The lab uses BIND (named) as its authoritative DNS server. The DNS infrastructure is managed by the \texttt{dns} Ansible role in the lab-franklin collection, located at:
\begin{lstlisting}[style=mystyle]
ansible/collections/ansible_collections/lab/franklin/roles/dns/
    \end{lstlisting}

\subsection{Zones Maintained}

\begin{itemize}
    \item \textbf{Forward zone}: db.home.lab -- maps hostnames to IPs for all lab infrastructure
    
    \item \textbf{Reverse zones}: PTR records for each subnet (10.10.x.0/24 ranges)
    
    \item \textbf{SRV records}: Kerberos service discovery (\_kerberos.\_tcp, \_kdc.tcp)
\end{itemize}

\subsection{DNS Server Location}

The DNS server (ns1) is the single authoritative nameserver for the lab. It has been unreachable during recent infrastructure audits, which explains why many hosts cannot resolve hostnames -- ns1 \textit{is} the DNS server and it's offline.

When ns1 is available, it resolves all internal hostnames. Key entries include:
\begin{itemize}
    \item time.lab.bitsmasher.net $\rightarrow$ 10.10.12.2 (NTP server)
    \item odroid-c1.lab.bitsmasher.net $\rightarrow$ 10.10.12.254 (KDC)
    \item ldap.lab.bitsmasher.net $\rightarrow$ 10.10.13.1 (OpenLDAP)
    \item blowfish.lab.bitsmasher.net $\rightarrow$ 10.10.14.85 (reassigned from 10.10.12.15)
\end{itemize}

\subsection{Molecule Test Harness}

The DNS role has a molecule test harness that confirms:
\begin{enumerate}
    \item \textbf{Converge}: BIND packages installed, /etc/bind directory created, named.conf.options written
    \item \textbf{Test sequence}: dependency $\rightarrow$ syntax $\rightarrow$ create (podman container) $\rightarrow$ prepare $\rightarrow$ converge
    
    \item \textbf{Limitation}: molecule test uses Ubuntu 20.04 container with hostname "server1"; the role has a hostname gate that targets the real DNS server's hostname
\end{enumerate}

The DNS molecule test confirms correct package installation and configuration file placement but does not verify actual name resolution (no live DNS server in the container test environment).

\subsection{DNS Failure Impact}

When ns1 is offline:
\begin{itemize}
    \item Lab hosts fall back to /etc/hosts entries (which are authoritative but not dynamic)
    
    \item Kerberos KDC discovery via SRV records fails -- clients need manual kdc = entry in krb5.conf
    
    \item DHCP reservations on the Dream Machine remain functional (DHCP server is separate from DNS)
\end{itemize}

The DNS role's molecule test harness ensures that when the server comes back online, its configuration will install and write correctly.


\chapter{Core Services}

\section{Network Time (NTP)}
A two-tier NTP architecture provides time synchronization from the Stratum-1 GPS-referenced time host to all lab clients.

% \section{Network Time (NTP)}

\subsection{Overview}

Time synchronization is one of the most critical and least-visible aspects of network management. All lab services -- Kerberos authentication, TLS certificate validation, log correlation, distributed databases, and filesystem consistency -- depend on time accuracy. The lab uses a two-tier NTP architecture:

\begin{enumerate}
    \item \textbf{Stratum 1 source}: GPS-referenced clock on the \texttt{time} host (10.10.12.2) provides the lab's authoritative time reference.
    
    \item \textbf{Stratum 2+ clients}: All other hosts sync to \texttt{time.lab.bitsmasher.net}.
\end{enumerate}

\subsection{The Time Host}

\begin{itemize}
    \item \textbf{Hostname}: time.lab.bitsmasher.net (alias: time)
    
    \item \textbf{IP}: 10.10.12.2
    
    \item \textbf{User}: openclaw
    
    \item \textbf{OS}: Debian 12 bookworm
    
    \item \textbf{Role}: Stratum-1 NTP server / reference clock source for the lab
\end{itemize}

The time host has a GPS unit connected (serial/tty interface) providing UTC discipline. Its ntpd is configured via ntpsec to serve the lab subnet and use the GPS PPS signal as the primary time reference.

\subsection{Configuration on Client Hosts}

The NTP configuration is managed by Ansible via the \texttt{ntp} role in the lab-franklin collection. Here is what a typical client configuration looks like (as seen on stargate):

\begin{lstlisting}[style=mystyle]
# /etc/ntpsec/ntp.conf -- managed by Ansible
driftfile /var/lib/ntp/ntp.drift
leapfile /usr/share/zoneinfo/leap-seconds.list

statsdir /var/log/ntpstats/

statistics loopstats peerstats clockstats
filegen loopstats file loopstats type day enable
filegen peerstats file peerstats type day enable
filegen clockstats file clockstats type day enable

server time.lab.bitsmasher.net
restrict 10.10.8.0/21
restrict -4 default
\end{lstlisting}

Key points:
\begin{itemize}
    \item \texttt{driftfile}: Tracks oscillator drift between NTP restarts. Critical for stability on hosts that reboot frequently (Jetson devices, test VMs).
    
    \item \texttt{statistics}: Logs loop/peer/clock stats daily in /var/log/ntpstats/. Use these to verify sync health.
    
    \item \texttt{server}: Points to the lab's time host -- never pool.ntp.org for infrastructure hosts that need deterministic time sources.
    
    \item \texttt{restrict}: 
    \begin{itemize}
        \item \texttt{10.10.8.0/21} grants query access to the main lab subnet (used by other internal clients or monitoring tools).
        
        \item \texttt{-4 default} blocks all IPv4 queries not explicitly allowed, providing a deny-by-default posture.
    \end{itemize}
\end{itemize}

\subsection{Installation and Setup Steps}

For a new host that needs NTP:

\begin{enumerate}
    \item Install ntpsec packages:
    \begin{lstlisting}[style=mystyle]
apt update && apt install -y ntpsec
    \end{lstlisting}

    \item Configure /etc/ntpsec/ntp.conf (either manually or via Ansible ntp role):
    \begin{enumerate}
        \item Set \texttt{server time.lab.bitsmasher.net} as the upstream reference.
        
        \item Set timezone: \texttt{timedatectl set-timezone America/Denver} (or appropriate TZ).
        
        \item Create the log directory if it doesn't exist: \texttt{mkdir -p /var/log/ntpstats}.
    \end{enumerate}

    \item Start and enable the service:
    \begin{lstlisting}[style=mystyle]
systemctl enable --now ntpsec.service
timedatectl  # verify "System clock synchronized: yes"
    \end{lstlisting}
\end{enumerate}

\subsection{Verification on Clients}

Check synchronization status with these commands:

\begin{lstlisting}[style=mystyle]
# Check if the system clock is synchronized
timedatectl | grep "System clock synchronized"

# Check NTP peer status (if ntpq is available)
ntpq -p

# View recent sync logs
journalctl -u ntpsec --since today
    \end{lstlisting}

\subsection{Notes for Deployment}

\begin{itemize}
    \item The Ansible \texttt{ntp} role handles all of the above automatically across lab hosts. It installs packages, copies configuration, sets timezone, and creates log directories.
    
    \item For Jetson devices (node900--node903), drift compensation is especially important -- they boot frequently and have less stable oscillators than desktop/server hardware.
    
    \item The GPS unit on the time host should be verified periodically: confirm it's locking to satellites, not using its internal oscillator as a holdover source. Check this by running \texttt{ntptime} on the time host itself.
    
    \item If the time host goes offline and all clients lose sync beyond their maximum correction threshold (typically 128 seconds), systems may need manual time correction before Kerberos/TLS will function again.
\end{itemize}

\section{Network Time (NTP)}

\subsection{Overview}

Time synchronization is one of the most critical and least-visible aspects of network management. All lab services -- Kerberos authentication, TLS certificate validation, log correlation, distributed databases, and filesystem consistency -- depend on time accuracy. The lab uses a two-tier NTP architecture:

\begin{enumerate}
    \item \textbf{Stratum 1 source}: GPS-referenced clock on the \texttt{time} host (10.10.12.2) provides the lab's authoritative time reference.
    
    \item \textbf{Stratum 2+ clients}: All other hosts sync to \texttt{time.lab.bitsmasher.net}.
\end{enumerate}

\subsection{The Time Host}

\begin{itemize}
    \item \textbf{Hostname}: time.lab.bitsmasher.net (alias: time)
    
    \item \textbf{IP}: 10.10.12.2
    
    \item \textbf{User}: openclaw
    
    \item \textbf{OS}: Debian 12 bookworm
    
    \item \textbf{Role}: Stratum-1 NTP server / reference clock source for the lab
\end{itemize}

The time host has a GPS unit connected (serial/tty interface) providing UTC discipline. Its ntpd is configured via ntpsec to serve the lab subnet and use the GPS PPS signal as the primary time reference.

\subsection{Configuration on Client Hosts}

The NTP configuration is managed by Ansible via the \texttt{ntp} role in the lab-franklin collection. Here is what a typical client configuration looks like (as seen on stargate):

\begin{lstlisting}[style=mystyle]
# /etc/ntpsec/ntp.conf -- managed by Ansible
driftfile /var/lib/ntp/ntp.drift
leapfile /usr/share/zoneinfo/leap-seconds.list

statsdir /var/log/ntpstats/

statistics loopstats peerstats clockstats
filegen loopstats file loopstats type day enable
filegen peerstats file peerstats type day enable
filegen clockstats file clockstats type day enable

server time.lab.bitsmasher.net
restrict 10.10.8.0/21
restrict -4 default
\end{lstlisting}

Key points:
\begin{itemize}
    \item \texttt{driftfile}: Tracks oscillator drift between NTP restarts. Critical for stability on hosts that reboot frequently (Jetson devices, test VMs).
    
    \item \texttt{statistics}: Logs loop/peer/clock stats daily in /var/log/ntpstats/. Use these to verify sync health.
    
    \item \texttt{server}: Points to the lab's time host -- never pool.ntp.org for infrastructure hosts that need deterministic time sources.
    
    \item \texttt{restrict}: 
    \begin{itemize}
        \item \texttt{10.10.8.0/21} grants query access to the main lab subnet (used by other internal clients or monitoring tools).
        
        \item \texttt{-4 default} blocks all IPv4 queries not explicitly allowed, providing a deny-by-default posture.
    \end{itemize}
\end{itemize}

\subsection{Installation and Setup Steps}

For a new host that needs NTP:

\begin{enumerate}
    \item Install ntpsec packages:
    \begin{lstlisting}[style=mystyle]
apt update && apt install -y ntpsec
    \end{lstlisting}

    \item Configure /etc/ntpsec/ntp.conf (either manually or via Ansible ntp role):
    \begin{enumerate}
        \item Set \texttt{server time.lab.bitsmasher.net} as the upstream reference.
        
        \item Set timezone: \texttt{timedatectl set-timezone America/Denver} (or appropriate TZ).
        
        \item Create the log directory if it doesn't exist: \texttt{mkdir -p /var/log/ntpstats}.
    \end{enumerate}

    \item Start and enable the service:
    \begin{lstlisting}[style=mystyle]
systemctl enable --now ntpsec.service
timedatectl  # verify "System clock synchronized: yes"
    \end{lstlisting}
\end{enumerate}

\subsection{Verification on Clients}

Check synchronization status with these commands:

\begin{lstlisting}[style=mystyle]
# Check if the system clock is synchronized
timedatectl | grep "System clock synchronized"

# Check NTP peer status (if ntpq is available)
ntpq -p

# View recent sync logs
journalctl -u ntpsec --since today
    \end{lstlisting}

\subsection{Notes for Deployment}

\begin{itemize}
    \item The Ansible \texttt{ntp} role handles all of the above automatically across lab hosts. It installs packages, copies configuration, sets timezone, and creates log directories.
    
    \item For Jetson devices (node900--node903), drift compensation is especially important -- they boot frequently and have less stable oscillators than desktop/server hardware.
    
    \item The GPS unit on the time host should be verified periodically: confirm it's locking to satellites, not using its internal oscillator as a holdover source. Check this by running \texttt{ntptime} on the time host itself.
    
    \item If the time host goes offline and all clients lose sync beyond their maximum correction threshold (typically 128 seconds), systems may need manual time correction before Kerberos/TLS will function again.
\end{itemize}


\section{Kerberos Authentication}
The KDC on odroid-c1 runs MIT Kerberos for the lab.bitsmasher.net realm, managing authentication for all services.

% \section{Kerberos Authentication}

\subsection{Realm Configuration}

\begin{itemize}
    \item \textbf{Realm name}: lab.bitsmasher.net
    
    \item \textbf{KDC host}: odroid-c1.lab.bitsmasher.net (10.10.12.254)
    
    \item \textbf{KDC software}: MIT Kerberos KDC running on odroid-c1
    
    \item \textbf{Uptime}: ~242 days (stable, rarely rebooted)
\end{itemize}

\subsection{Access and Key Management}

The KDC only has SSH key auth configured for the \texttt{root} user. The \texttt{franklin} user has been rejected by the KDC's SSH configuration -- direct root access is required for any management tasks.

Key file: \texttt{\textasciitilde/.ssh/id\_ed25519\_openclaw} (the standard openclaw key).

\subsection{Principal Conventions}

Kerberos principals in the lab.bitsmasher.net realm should follow these conventions:
\begin{itemize}
    \item Service principals: \texttt{service/FQDN@lab.BITSMASHER.NET}
    \item User principals: \texttt{username@lab.BITSMASHER.NET}
    \item Host principals: \texttt{host/FQDN@lab.BITSMASHER.NET}
\end{itemize}

\subsection{DNS/Kerberos Integration -- RFC Compliant (August 2026)}

Kerberos relies on DNS SRV records for KDC discovery. The following records point directly to FQDN A records (no CNAME indirection):

\begin{lstlisting}[style=mystyle]
_kerberos-adm._tcp.lab.bitsmasher.net IN SRV 0 100 749 kdc1.lab.bitsmasher.net.
_kerberos._tcp.lab.bitsmasher.net     IN SRV 0 0 88  kdc1.lab.bitsmasher.net.
_kerberos._udp.lab.bitsmasher.net     IN SRV 0 0 88  kdc1.lab.bitsmasher.net.
_kdc._tcp.lab.bitsmasher.net          IN SRV 0 0 88  kdc1.lab.bitsmasher.net.
_kdc._udp.lab.bitsmasher.net          IN SRV 0 0 88  kdc1.lab.bitsmasher.net.
    \end{lstlisting}

The \texttt{_kerberos-adm.\_tcp} record now points directly to the A record for \texttt{kdc1.lab.bitsmasher.net}, per RFC 2782 compliance.

These records are authoritative on node3 (10.10.12.3), the new DNS master following the server1 decommissioning. If DNS is unavailable, clients fall back to manual \texttt{kdc = kdc1.lab.bitsmasher.net} in \texttt{/etc/krb5.conf}.

\subsection{Keytab Management}

Keytabs for service principals should be:
\begin{itemize}
    \item Generated on the KDC host (odroid-c1) using \texttt{kadmin.local} or \texttt{kadmin -q}
    \item Transferred securely to target hosts (via scp with SSH keys)
    \item Set to mode 600 and owned by the appropriate service user
\end{itemize}

\textbf{Note}: Kerberos keytabs have not been deployed across the infrastructure. NFS exports with \texttt{sec=krb5i} are currently non-functional pending keytab deployment. The Ansible ssh role handles the StrictModes requirements for SSH-based keytab transfer.

\section{Kerberos Authentication}

\subsection{Realm Configuration}

\begin{itemize}
    \item \textbf{Realm name}: lab.bitsmasher.net
    
    \item \textbf{KDC host}: odroid-c1.lab.bitsmasher.net (10.10.12.254)
    
    \item \textbf{KDC software}: MIT Kerberos KDC running on odroid-c1
    
    \item \textbf{Uptime}: ~242 days (stable, rarely rebooted)
\end{itemize}

\subsection{Access and Key Management}

The KDC only has SSH key auth configured for the \texttt{root} user. The \texttt{franklin} user has been rejected by the KDC's SSH configuration -- direct root access is required for any management tasks.

Key file: \texttt{\textasciitilde/.ssh/id\_ed25519\_openclaw} (the standard openclaw key).

\subsection{Principal Conventions}

Kerberos principals in the lab.bitsmasher.net realm should follow these conventions:
\begin{itemize}
    \item Service principals: \texttt{service/FQDN@lab.BITSMASHER.NET}
    \item User principals: \texttt{username@lab.BITSMASHER.NET}
    \item Host principals: \texttt{host/FQDN@lab.BITSMASHER.NET}
\end{itemize}

\subsection{DNS/Kerberos Integration -- RFC Compliant (August 2026)}

Kerberos relies on DNS SRV records for KDC discovery. The following records point directly to FQDN A records (no CNAME indirection):

\begin{lstlisting}[style=mystyle]
_kerberos-adm._tcp.lab.bitsmasher.net IN SRV 0 100 749 kdc1.lab.bitsmasher.net.
_kerberos._tcp.lab.bitsmasher.net     IN SRV 0 0 88  kdc1.lab.bitsmasher.net.
_kerberos._udp.lab.bitsmasher.net     IN SRV 0 0 88  kdc1.lab.bitsmasher.net.
_kdc._tcp.lab.bitsmasher.net          IN SRV 0 0 88  kdc1.lab.bitsmasher.net.
_kdc._udp.lab.bitsmasher.net          IN SRV 0 0 88  kdc1.lab.bitsmasher.net.
    \end{lstlisting}

The \texttt{_kerberos-adm.\_tcp} record now points directly to the A record for \texttt{kdc1.lab.bitsmasher.net}, per RFC 2782 compliance.

These records are authoritative on node3 (10.10.12.3), the new DNS master following the server1 decommissioning. If DNS is unavailable, clients fall back to manual \texttt{kdc = kdc1.lab.bitsmasher.net} in \texttt{/etc/krb5.conf}.

\subsection{Keytab Management}

Keytabs for service principals should be:
\begin{itemize}
    \item Generated on the KDC host (odroid-c1) using \texttt{kadmin.local} or \texttt{kadmin -q}
    \item Transferred securely to target hosts (via scp with SSH keys)
    \item Set to mode 600 and owned by the appropriate service user
\end{itemize}

\textbf{Note}: Kerberos keytabs have not been deployed across the infrastructure. NFS exports with \texttt{sec=krb5i} are currently non-functional pending keytab deployment. The Ansible ssh role handles the StrictModes requirements for SSH-based keytab transfer.


\section{LDAP Directory Services}
OpenLDAP (slapd) on bbb1 provides directory-based identity management. Service has been offline since December 2025 due to a TLS certificate issue.

% \section{LDAP Directory Services}

\subsection{Server Configuration}

\begin{itemize}
    \item \textbf{Hostname}: ldap.lab.bitsmasher.net (alias: bbb1)
    
    \item \textbf{IP}: 10.10.13.1
    
    \item \textbf{Software}: OpenLDAP (slapd)
    
    \item \textbf{Status}: SSH reachable as root, but slapd service has been \textbf{failed since December 2025}
\end{itemize}

\subsection{Outage Details}

The slapd failure manifests as:
\begin{lstlisting}[style=mystyle]
TLS init def ctx failed (-1)
    \end{lstlisting}

This indicates a TLS certificate configuration issue -- likely an expired, missing, or misconfigured CA certificate. The server is listening on its network interface (SSH works for root), but the LDAP daemon itself has not recovered since December 2025.

\subsection{Monitoring Approach}

Franklin maintains \texttt{/home/franklin/status.sh} which checks:
\begin{itemize}
    \item slapd service status on bbb1/ldap
    
    \item ldapsearch queries for franklin and sly DN lookups -- both fail because the directory is down
\end{itemize}

The script provides a simple pass/fail indicator for LDAP health without requiring complex monitoring infrastructure.

\subsection{Recovery Steps}

To restore the LDAP server:
\begin{enumerate}
    \item SSH to ldap.lab.bitsmasher.net as root (ed25519 key)
    
    \item Check slapd logs: \texttt{journalctl -u slapd --since "2025-12-13"}
    
    \item Verify TLS certificate chain in /etc/ssl or the configured cert path
    
    \item Update or regenerate the CA-signed certificate if expired
    
    \item Restart slapd: \texttt{systemctl restart slapd}
    
    \item Verify with ldapsearch for franklin and sly DNs
\end{enumerate}

\subsection{Directory Structure (Expected)}

The OpenLDAP directory should contain entries for:
\begin{itemize}
    \item User entries for franklin, Sly (slyborg), and other lab members
    
    \item Group/organizational unit entries for lab access control
    
    \item Service accounts for LDAP-authorized services (Kerberos integration, web auth)
\end{itemize}

Once restored, the LDAP directory serves as the central authentication source for the lab, complementing Kerberos for identity management.

\section{LDAP Directory Services}

\subsection{Server Configuration}

\begin{itemize}
    \item \textbf{Hostname}: ldap.lab.bitsmasher.net (alias: bbb1)
    
    \item \textbf{IP}: 10.10.13.1
    
    \item \textbf{Software}: OpenLDAP (slapd)
    
    \item \textbf{Status}: SSH reachable as root, but slapd service has been \textbf{failed since December 2025}
\end{itemize}

\subsection{Outage Details}

The slapd failure manifests as:
\begin{lstlisting}[style=mystyle]
TLS init def ctx failed (-1)
    \end{lstlisting}

This indicates a TLS certificate configuration issue -- likely an expired, missing, or misconfigured CA certificate. The server is listening on its network interface (SSH works for root), but the LDAP daemon itself has not recovered since December 2025.

\subsection{Monitoring Approach}

Franklin maintains \texttt{/home/franklin/status.sh} which checks:
\begin{itemize}
    \item slapd service status on bbb1/ldap
    
    \item ldapsearch queries for franklin and sly DN lookups -- both fail because the directory is down
\end{itemize}

The script provides a simple pass/fail indicator for LDAP health without requiring complex monitoring infrastructure.

\subsection{Recovery Steps}

To restore the LDAP server:
\begin{enumerate}
    \item SSH to ldap.lab.bitsmasher.net as root (ed25519 key)
    
    \item Check slapd logs: \texttt{journalctl -u slapd --since "2025-12-13"}
    
    \item Verify TLS certificate chain in /etc/ssl or the configured cert path
    
    \item Update or regenerate the CA-signed certificate if expired
    
    \item Restart slapd: \texttt{systemctl restart slapd}
    
    \item Verify with ldapsearch for franklin and sly DNs
\end{enumerate}

\subsection{Directory Structure (Expected)}

The OpenLDAP directory should contain entries for:
\begin{itemize}
    \item User entries for franklin, Sly (slyborg), and other lab members
    
    \item Group/organizational unit entries for lab access control
    
    \item Service accounts for LDAP-authorized services (Kerberos integration, web auth)
\end{itemize}

Once restored, the LDAP directory serves as the central authentication source for the lab, complementing Kerberos for identity management.


\chapter{Security}

\section{Authentication and Access Control}

SSH key management, password policy, and user accounts define access to the lab infrastructure. All automation hosts share a common ed25519 key pair for the openclaw user. The primary admin (franklin) has sudo NOPASSWD on skynet; root is restricted to KDC, ldap, and music host only.

\section{Network Security}

Subnet segmentation provides implicit isolation across the lab. The ``prereq'' Ansible role handles firewall exceptions for K3s API ports on each host (6443, 2379--2381). Both UFW (Debian) and Firewalld (RedHat-family) are supported.

\section{Certificate and TLS Management}

The ``tls'' role manages certificate provisioning across the lab. The OpenLDAP server experienced a TLS failure in December 2025 due to an expired certificate -- this remains the primary ongoing TLS vulnerability.

\section{Credential Management}

Credentials use pass (password manager) with GPG encryption. The Stash House project manages encrypted credential backups synced via bare git repos. Sensitive Ansible variables are stored in ansible-vault files that should never be committed to version control.

\section{Backup and Disaster Recovery}

NFS-mounted clusterfs2 provides shared storage with redundant copies across hosts. DNS zones are version-controlled in BIND configuration files.

Service availability:

\begin{itemize}
  \item \textbf{NTP stratum-1} (time.lab.bitsmasher.net) -- Online
  \item \textbf{Kerberos KDC} (odroid-c1.lab.bitsmasher.net) -- Online, ~242 days uptime
  \item \textbf{DNS BIND} (ns1.lab.bitsmasher.net) -- Offline since December 2025
  \item \textbf{OpenLDAP} (ldap.lab.bitsmasher.net) -- Offline since December 2025
  \item \textbf{Minecraft} (skynet.lab.bitsmasher.net) -- Online
\end{itemize}

Disaster recovery priority: NTP > DNS > Kerberos KDC > LDAP. Time sync failure cascades to all Kerberos authentication; DNS failure makes hosts unreachable by name.

\section{Container and Kubernetes Security}

The k3s cluster is hardened with version pinning (2.1.x), Containerd runtime with NVIDIA GPU support, TLS for all API communication (k3s default), and pod security policies via Kubernetes native admission controllers.

\section{Automated Security Scanning}

The .github directory contains these automated workflows: bandit.yml (Bandit Python security), tfsec-sarif.yml (Terraform IaC security), trivy.yaml + trivy-cb.yml (container/filesystem vulnerability scanning), bash\_check.yaml (GNU autotools build validation with ``make security''), and markdown.yml (markdownlint). All run on every pull request and main branch push.

\chapter{Storage}

\section{NFS Mounts and Backup}
clusterfs2 provides shared NFS storage across hosts. Credential backups use GPG-encrypted pass entries synced via bare git repos (Stash House project).

% \section{Storage Infrastructure}

\subsection{NFS Mounts}

The lab relies on NFS for shared storage across hosts. Key mount points:

\begin{itemize}
    \item \textbf{/mnt/clusterfs2}: Primary cluster workspace mounted on chonk and stargate (from the NFS server). This is where Ansible roles, molecule tests, and project files live.
    
    \item \textbf{/mnt/snowy}: Formerly a full NFS share; now just a mounted drive after the hard disk was physically removed. Should be audited for stale mount entries.
\end{itemize}

The NFS role in lab-franklin's Ansible collection handles:
\begin{itemize}
    \item Server-side export configuration (/etc/exports)
    \item Client-side fstab management and mount points
    \item Mount path conventions per role (previously /mnt/clusterfs, /mnt/backup1, /mnt/storage{1,2,3})
\end{itemize}

Note: The NFS role has hostname gates (\texttt{'snowy' in inventory\_hostname}) that prevent testing on molecule containers. Hardcoded mount paths depend on real NFS servers being online. This is a known limitation of the test harness.

\subsection{Snapshot and Backup Strategy}

The lab uses bare git repos for version control and backup:
\begin{itemize}
    \item GPG-encrypted repositories synced across hosts for credential protection (Stash House project)
    \item Pass-based password store with GPG backend
    \item No LFS -- ABSOLUTE RULE: do not use git-lfs anywhere. Microsoft once held Minecraft chunk files hostage when a repo hit the 10GB GitHub limit with LFS enabled.
\end{itemize}

\subsection{Local Storage Considerations}

\begin{itemize}
    \item Jetson nodes (node90x): limited eMMC storage; rely on NFS for workspace data
    \item chonk: large local disk, primary build and gateway host
    \item stargate: Debian 12 box with full Ansible workspace mirror
\end{itemize}

\subsection{LFS Policy Reminder}

\textbf{ABSOLUTE RULE: Do NOT use git-lfs anywhere in any lab-franklin repository.}

This applies across all repos: training-ai, minecraft, Ansible collections, everything. The incident with Minecraft chunk files being held hostage by GitHub when the repo hit 10GB with LFS is a real precedent -- never risk it.

\section{Storage Infrastructure}

\subsection{NFS Mounts}

The lab relies on NFS for shared storage across hosts. Key mount points:

\begin{itemize}
    \item \textbf{/mnt/clusterfs2}: Primary cluster workspace mounted on chonk and stargate (from the NFS server). This is where Ansible roles, molecule tests, and project files live.
    
    \item \textbf{/mnt/snowy}: Formerly a full NFS share; now just a mounted drive after the hard disk was physically removed. Should be audited for stale mount entries.
\end{itemize}

The NFS role in lab-franklin's Ansible collection handles:
\begin{itemize}
    \item Server-side export configuration (/etc/exports)
    \item Client-side fstab management and mount points
    \item Mount path conventions per role (previously /mnt/clusterfs, /mnt/backup1, /mnt/storage{1,2,3})
\end{itemize}

Note: The NFS role has hostname gates (\texttt{'snowy' in inventory\_hostname}) that prevent testing on molecule containers. Hardcoded mount paths depend on real NFS servers being online. This is a known limitation of the test harness.

\subsection{Snapshot and Backup Strategy}

The lab uses bare git repos for version control and backup:
\begin{itemize}
    \item GPG-encrypted repositories synced across hosts for credential protection (Stash House project)
    \item Pass-based password store with GPG backend
    \item No LFS -- ABSOLUTE RULE: do not use git-lfs anywhere. Microsoft once held Minecraft chunk files hostage when a repo hit the 10GB GitHub limit with LFS enabled.
\end{itemize}

\subsection{Local Storage Considerations}

\begin{itemize}
    \item Jetson nodes (node90x): limited eMMC storage; rely on NFS for workspace data
    \item chonk: large local disk, primary build and gateway host
    \item stargate: Debian 12 box with full Ansible workspace mirror
\end{itemize}

\subsection{LFS Policy Reminder}

\textbf{ABSOLUTE RULE: Do NOT use git-lfs anywhere in any lab-franklin repository.}

This applies across all repos: training-ai, minecraft, Ansible collections, everything. The incident with Minecraft chunk files being held hostage by GitHub when the repo hit 10GB with LFS is a real precedent -- never risk it.


\section{LFS Policy}
Absolute rule: no git-lfs in any lab repository -- see the storage section for details on the incident that made this policy necessary.

\chapter{Internal PyPI Repository}

\section{Overview}
An internal PyPI server provides Python packages to development machines without depending on external PyPI.org servers. Packages are stored on NFS-mounted storage and served via HTTP from the web document root.

% \section{Internal PyPI Repository}

The lab runs an internal PyPI package repository to provide Python packages to development machines without depending on external PyPI servers. This avoids network bottlenecks, supports offline development, and enables use of proprietary or lab-specific packages not available on PyPI.org.

\subsection{Architecture}

The PyPI directory lives at \texttt{/mnt/storage1/LAB/pypi}. A symlink is created from the web server's document root (\texttt{/var/www/html/pypi}) pointing to this directory, making packages accessible via HTTP.

Access URL pattern:
\begin{lstlisting}[style=mystyle]
http://<storage-host>/pypi/simple/
    \end{lstlisting}

This follows the PEP 503 simple repository API structure expected by pip.

\subsection{Configuration via Ansible}

The \texttt{pypi\_internal} role in the lab/franklin collection automates this setup:

\begin{itemize}
    \item \textbf{Role path}: \texttt{roles/pypi\_internal/}
    
    \item \textbf{Variables}:
    \begin{itemize}
        \item \texttt{html\_dir}: web root (\texttt{/var/www/html})
        \item \texttt{pypi\_dir}: package directory (\texttt{/mnt/storage1/LAB/pypi})
    \end{itemize}
    
    \item \textbf{Task}: Creates a symbolic link from \texttt{\$html\_dir/pypi} $\rightarrow$ \texttt{\$pypi\_dir}
    
    \item \textbf{Playbook integration}: Included in the cluster/storage playbook alongside NFS, LDAP, Kerberos, DNS, and ntp roles
\end{itemize}

\subsection{Usage from Client Machines}

Configure pip to use the internal PyPI server:
\begin{lstlisting}[style=mystyle]
# In ~/.pip/pip.conf or /etc/pip.conf:
[global]
index-url = http://<storage-host>/pypi/simple/

# Optionally trust host if not using HTTPS
trusted-host = <storage-host>
    \end{lstlisting}

Or use pip's \texttt{\-\-index\-url} flag directly:
\begin{lstlisting}[style=mystyle]
pip install --index-url http://<storage-host>/pypi/simple/ some-package
    \end{lstlisting}

\subsection{Package Maintenance}

To add packages to the internal PyPI, place them in the pypi directory with PEP 503-compliant directory structure:
\begin{itemize}
    \item Package directories (lowercase names)
    
    \item File links inside each package directory pointing to the actual wheel/sdist files
    
    \item The web server must have read access to /mnt/storage1/LAB/pypi
\end{itemize}

Packages can be built and added using standard pip tools:
\begin{lstlisting}[style=mystyle]
pip install --no-index --find-links=/mnt/storage1/LAB/pypi/simple/ <package>
    \end{lstlisting}

\section{Internal PyPI Repository}

The lab runs an internal PyPI package repository to provide Python packages to development machines without depending on external PyPI servers. This avoids network bottlenecks, supports offline development, and enables use of proprietary or lab-specific packages not available on PyPI.org.

\subsection{Architecture}

The PyPI directory lives at \texttt{/mnt/storage1/LAB/pypi}. A symlink is created from the web server's document root (\texttt{/var/www/html/pypi}) pointing to this directory, making packages accessible via HTTP.

Access URL pattern:
\begin{lstlisting}[style=mystyle]
http://<storage-host>/pypi/simple/
    \end{lstlisting}

This follows the PEP 503 simple repository API structure expected by pip.

\subsection{Configuration via Ansible}

The \texttt{pypi\_internal} role in the lab/franklin collection automates this setup:

\begin{itemize}
    \item \textbf{Role path}: \texttt{roles/pypi\_internal/}
    
    \item \textbf{Variables}:
    \begin{itemize}
        \item \texttt{html\_dir}: web root (\texttt{/var/www/html})
        \item \texttt{pypi\_dir}: package directory (\texttt{/mnt/storage1/LAB/pypi})
    \end{itemize}
    
    \item \textbf{Task}: Creates a symbolic link from \texttt{\$html\_dir/pypi} $\rightarrow$ \texttt{\$pypi\_dir}
    
    \item \textbf{Playbook integration}: Included in the cluster/storage playbook alongside NFS, LDAP, Kerberos, DNS, and ntp roles
\end{itemize}

\subsection{Usage from Client Machines}

Configure pip to use the internal PyPI server:
\begin{lstlisting}[style=mystyle]
# In ~/.pip/pip.conf or /etc/pip.conf:
[global]
index-url = http://<storage-host>/pypi/simple/

# Optionally trust host if not using HTTPS
trusted-host = <storage-host>
    \end{lstlisting}

Or use pip's \texttt{\-\-index\-url} flag directly:
\begin{lstlisting}[style=mystyle]
pip install --index-url http://<storage-host>/pypi/simple/ some-package
    \end{lstlisting}

\subsection{Package Maintenance}

To add packages to the internal PyPI, place them in the pypi directory with PEP 503-compliant directory structure:
\begin{itemize}
    \item Package directories (lowercase names)
    
    \item File links inside each package directory pointing to the actual wheel/sdist files
    
    \item The web server must have read access to /mnt/storage1/LAB/pypi
\end{itemize}

Packages can be built and added using standard pip tools:
\begin{lstlisting}[style=mystyle]
pip install --no-index --find-links=/mnt/storage1/LAB/pypi/simple/ <package>
    \end{lstlisting}


\chapter{Container Orchestration}

\section{k3s Cluster}
k3s provides Kubernetes across the infrastructure with version pinning (2.1.x), containerd runtime, and NVIDIA Container Toolkit for GPU workloads.

% \section{Kubernetes (k3s) Cluster}

\subsection{Architecture}

The lab uses k3s (Lightweight Kubernetes from Rancher) for container orchestration across the infrastructure:

\begin{itemize}
    \item \textbf{k3s\_server}: Master node running k3s control plane (host: head2, referenced but unreachable in current audit)
    
    \item \textbf{k3s\_agent}: Worker nodes that register with the server and receive workloads
    
    \item \textbf{GPU nodes}: Specialized workers for NVIDIA container runtime and compute workloads
\end{itemize}

\subsection{Version Pinning}

k3s version is pinned to the 2.1.x series (specifically 2.1.5 as of the latest network\_update.sh). This prevents surprise breakage from upstream k3s releases. The pinning policy follows:
\begin{lstlisting}[style=mystyle]
K3S_VERSION="2.1.5"
curl -sfL https://get.k3s.io | sh -s - \
    --write-kubeconfig-mode 644
    \end{lstlisting}

\subsection{Container Runtime}

k3s uses containerd as the default runtime, with NVIDIA Container Toolkit for GPU workloads:
\begin{itemize}
    \item \textbf{nvidia-container-toolkit}: Installed on GPU nodes via the network\_update.sh automation
    
    \item \textbf{containerd config}: Located at \texttt{/var/lib/rancher/k3s/agent/etc/containerd/config.toml}
    
    \item NVIDIA runtime configured in containerd to expose GPU devices to containers
\end{itemize}

\subsection{Cluster Topology}

\begin{tabularx}{\textwidth}{l l X}
\textbf{Node Role} & \textbf{Target Hosts} & \textbf{Notes} \\
\hline
k3s\_server & head2 (planned) & Control plane; not yet online\\
k3s\_agent & node900, node901, node903 & Jetson nodes as workers\\
nvidia\_nodes & GPU-equipped hosts & Container Toolkit + k3s agent
\end{tabularx}

\subsection{kubeconfig Access}

After installation, kubeconfig is written to \texttt{/etc/rancher/k3s/k3s.yaml} with mode 644. Cluster access requires:
\begin{itemize}
    \item SSH key auth to the k3s\_server host
    
    \item Reading the kubeconfig file (or copying it securely)
    
    \item \texttt{kubectl --kubeconfig /etc/rancher/k3s/k3s.yaml cluster-info} to verify connectivity
\end{itemize}

The k3s_server role and k3s_agent role in lab-franklin handle the installation and configuration of these nodes via Ansible.

\section{Kubernetes (k3s) Cluster}

\subsection{Architecture}

The lab uses k3s (Lightweight Kubernetes from Rancher) for container orchestration across the infrastructure:

\begin{itemize}
    \item \textbf{k3s\_server}: Master node running k3s control plane (host: head2, referenced but unreachable in current audit)
    
    \item \textbf{k3s\_agent}: Worker nodes that register with the server and receive workloads
    
    \item \textbf{GPU nodes}: Specialized workers for NVIDIA container runtime and compute workloads
\end{itemize}

\subsection{Version Pinning}

k3s version is pinned to the 2.1.x series (specifically 2.1.5 as of the latest network\_update.sh). This prevents surprise breakage from upstream k3s releases. The pinning policy follows:
\begin{lstlisting}[style=mystyle]
K3S_VERSION="2.1.5"
curl -sfL https://get.k3s.io | sh -s - \
    --write-kubeconfig-mode 644
    \end{lstlisting}

\subsection{Container Runtime}

k3s uses containerd as the default runtime, with NVIDIA Container Toolkit for GPU workloads:
\begin{itemize}
    \item \textbf{nvidia-container-toolkit}: Installed on GPU nodes via the network\_update.sh automation
    
    \item \textbf{containerd config}: Located at \texttt{/var/lib/rancher/k3s/agent/etc/containerd/config.toml}
    
    \item NVIDIA runtime configured in containerd to expose GPU devices to containers
\end{itemize}

\subsection{Cluster Topology}

\begin{tabularx}{\textwidth}{l l X}
\textbf{Node Role} & \textbf{Target Hosts} & \textbf{Notes} \\
\hline
k3s\_server & head2 (planned) & Control plane; not yet online\\
k3s\_agent & node900, node901, node903 & Jetson nodes as workers\\
nvidia\_nodes & GPU-equipped hosts & Container Toolkit + k3s agent
\end{tabularx}

\subsection{kubeconfig Access}

After installation, kubeconfig is written to \texttt{/etc/rancher/k3s/k3s.yaml} with mode 644. Cluster access requires:
\begin{itemize}
    \item SSH key auth to the k3s\_server host
    
    \item Reading the kubeconfig file (or copying it securely)
    
    \item \texttt{kubectl --kubeconfig /etc/rancher/k3s/k3s.yaml cluster-info} to verify connectivity
\end{itemize}

The k3s_server role and k3s_agent role in lab-franklin handle the installation and configuration of these nodes via Ansible.


\chapter{Playbook Automation}

% \section{Ansible Automation}

\subsection{Current Landscape (August 2026)}

As of August 2026, the Ansible ecosystem is split between two distribution paths:

\begin{itemize}
    \item \textbf{Ansible community package} -- the traditional monolithic install (pip install ansible). Current version is 13.x (built on ansible-core 2.20). This includes 85+ collections with thousands of modules and plugins bundled together. Follows semantic versioning, releases two major versions per year plus monthly minor releases.
    
    \item \textbf{ansible-core} -- the minimal core engine. Maintains three versions at a time (current + two older). Versioned independently from the community package (2.19, 2.20, 2.21). Releases every four weeks for minor updates.

\end{itemize}

\textbf{EOL versions as of August 2026}: Ansible 10--12 and ansible-core 2.15--2.18 are EOL. Do not use these in new work. Only Ansible 13.x / ansible-core 2.20 is current; 14.0.0/2.21 is in development.

\subsection{Molecule Testing -- Status and Deprecation}

Molecule (the canonical integration testing framework for Ansible roles and playbooks) remains \textbf{active} as of August 2026 but its role has shifted:

\begin{itemize}
    \item Molecule v26.4.x is the latest series, released June 2026. It still supports only the latest two major versions of Ansible (N/N-1).
    
    \item Molecule is being superseded as the preferred test strategy by:
    \begin{itemize}
        \item \texttt{pytest-ansible} -- pytest extension for testing Ansible module/plugin Python code directly.
        
        \item \texttt{tox-ansible} -- integration with tox to run tests across multiple Python interpreters and ansible-core versions simultaneously.
        
        \item Execution Environments (containers) as the primary test target, tested via \texttt{ansible-navigator}.
    \end{itemize}
    
    \item The \textbf{community.molecule} collection still exists but new projects should prefer pytest-ansible + tox-ansible for role/test harness work. Molecule is not yet officially deprecated -- it is maintained on a best-effort basis while the community transitions.
    
    \item For existing molecule test harnesses (such as those in lab-franklin's Ansible collection at \texttt{ansible/collections/ansible\_collections/lab/franklin/roles/}), they will continue to work with ansible-core 2.20 and pytest-6.x but should be audited for:
    \begin{itemize}
        \item Podman driver compatibility (molecule continues using podman or docker as default driver)
        \item Python version requirements (molecule requires Python $\ge$3.9)
        \item Plugin deprecations in newer collections
    \end{itemize}

\end{itemize}

\textbf{Note}: For lab-franklin's molecule tests, keep testing with podman for now. The container-based approach is still valid; only the \textit{preferred tooling around it} is shifting toward pytest-ansible + execution environments.

\chapter{lab/franklin Collection}

The lab/franklin Ansible collection lives at:
\begin{lstlisting}[style=mystyle]
/mnt/clusterfs2/workspace/lab-franklin/ansible/collections/ansible_collections/lab/franklin/
    \end{lstlisting}

It is the central automation artifact for the entire lab.bitsmasher.net infrastructure, managing provisioning, configuration, and testing across all hosts.

\section{Collection Manifest (galaxy.yml)}

The collection follows standard Ansible Galaxy conventions with namespace \texttt{lab} and name \texttt{franklin}:

\begin{itemize}
    \item \textbf{Version}: 1.0.0
    
    \item \textbf{License}: MIT
    
    \item \textbf{Author}: franklin <franklin@bitsmasher.net>
    
    \item \textbf{Repository}: \url{https://github.com/devsecfranklin/lab-franklin}
    
    \item \textbf{Documentation}: \url{https://github.com/devsecfranklin/lab-franklin/docs}
\end{itemize}

\section{Role Inventory}

The collection contains roles organized under the roles/ directory. As of August 2026, the following roles are maintained with molecule test harnesses:

\begin{lstlisting}[style=mystyle]
roles/
├── apt_mirror/          # APT mirror configuration
├── beagleboard/         # BeagleBoard device provisioning
├── common/              # Shared base configuration (bootstrap.sh, common utils)
├── container_registry/  # Docker/container registry setup
├── ctfd/                # CTFd platform deployment
├── desktop/             # Desktop environment provisioning
├── docker/              # Docker engine installation
├── documentation/       # Documentation generation roles
├── dhcp/                # DHCP server configuration (dhcpd.conf)
├── dns/                 # BIND/named DNS server
├── golang/              # Go language toolchain
├── jetson-nano/         # NVIDIA Jetson Nano provisioning
├── k3s_agent/           # k3s worker node
├── k3s_server/          # k3s control plane server
├── k8s/                 # Kubernetes generic setup
├── latex/               # LaTeX environment
├── logging/             # Centralized logging
├── music/               # Music service configuration
├── nfs/                 # NFS server/client
├── nix/                 # Nix package manager
├── openbsd/             # OpenBSD host bootstrap (blowfish target)
│   ├── tasks/main.yml   # Includes files.yml, timezone setup
│   └── files/           # bootstrap.sh, login.conf, hosts entries
├── paloalto/            # Palo Alto firewall config
├── prereq/              # Prerequisite packages
├── pypi_internal/       # Internal PyPI repository setup
│   ├── tasks/main.yml   # Symlink pypi_dir -> html_dir/pypi
│   └── vars/main.yml    # html_dir=/var/www/html, pypi_dir=/mnt/storage1/LAB/pypi
├── python/              # Python environment
├── raspberrypi/         # Raspberry Pi provisioning
├── samba/               # Samba file sharing
├── security/            # Hardening and security baselines
├── shell/               # Shell configuration
├── ssh/                 # SSH server hardening
├── tls/                 # TLS certificate management
├── chonk/               # chonk-specific configuration
├── cluster/             # Cluster-wide settings
    \end{lstlisting}

Roles currently \textbf{missing molecule tests}: edge, mail, odroid, website (incomplete roles). These have tasks/main.yml but no test harness.

\section{Workspace Structure}

The full project layout:
\begin{lstlisting}[style=mystyle]
/workspace/lab-franklin/
├── Makefile.am                  # Top-level autotools target (Python, Docker, dev)
├── configure.ac                 # Autotools config (autoconf)
├── bootstrap.sh                 # Cross-platform bootstrapper
├── network_update.sh            # One-shot maintenance script (hardened 2026-08-02)
├── ansible/                     # ANSIBLE_HOME root
│   ├── playbook.yml             # Main playbook
│   ├── hosts                    # Inventory file
│   └── collections/             # Collections path
│       ├── ansible_collections/ # Galaxy namespace
│           └── lab/franklin/    # The lab/franklin collection
│               ├── galaxy.yml
│               ├── roles/       # All role modules (37 total)
│               └── docs/latex/  # LaTeX docs build (autotools)
├── container/                   # Docker/Podman configs
├── terraform/                   # Terraform infrastructure definitions
├── bin/                         # Shared utility scripts (common.sh)
└── docs/manual/                 # Lab manual documentation (this document)
    \end{lstlisting}

\section{Build System (Autotools)}

The project uses GNU Autotools for the root-level build:

\begin{itemize}
    \item \textbf{configure.ac}: Defines package (lab-franklin v0.2), libtool support, Python $\ge$ 3.9 requirement, podman/ansible/gpg checks, git version tagging
    
    \item \textbf{Makefile.am}: Top-level targets include test (ansible-lint + molecule), security (venv bootstrap), dev (Python venv in BUILDDIR)
    
    \item \textbf{bootstrap.sh}: Cross-platform bootstrapper that detects OS (Debian/RedHat/OpenBSD/macOS/Linux) and installs platform-specific packages, then runs aclocal $\rightarrow$ autoreconf $\rightarrow$ automake $\rightarrow$ configure
    
    \item \textbf{docs/manual/Makefile.am}: Standalone autotools stub for LaTeX docs build (clean target removes *.aux files)
\end{itemize}

The \texttt{network\_update.sh} script was hardened on 2026-08-02 with: set -euo pipefail, auto-resolved paths, pinned K3s version validation, ANSIBLE_HOME environment checks, removed dangerous clush/apt-get blast radius, and dead code removal. It runs the main playbook at \texttt{ansible/playbook.yml} as its primary action.

\chapter{Testing from Stargate}

\section{Molecule Test Harnesses on stargate.lab.bitsmasher.net}

The stargate host (10.10.16.66, Debian 12 bookworm) serves as the primary testing platform for the lab/franklin Ansible collection. Testing is run from \texttt{/mnt/clusterfs2/workspace/lab-franklin/ansible/collections/ansible\_collections/lab/franklin/roles/\textless role\textgreater /molecule/}.

\subsection{Test Environment}

Molecule tests on stargate use:
\begin{itemize}
    \item \textbf{Driver}: podman (containers, not Docker -- avoids daemon dependency)
    
    \item \textbf{Base image}: Ubuntu 20.04 for most roles, with role-specific variations
    
    \item \textbf{Execution user}: openclaw user (has sudo -n, no password required)
    
    \item \textbf{Python}: $\ge$ 3.9 (required by both ansible and molecule)
\end{itemize}

\subsection{Test Execution Flow}

For a given role, the molecule test cycle is:
\begin{enumerate}
    \item \textbf{dependency}: Install collection dependencies from galaxy.yml
    
    \item \textbf{syntax}: Validate playbook YAML syntax
    
    \item \textbf{create}: Launch podman container (e.g., "server1") as the test target host
    
    \item \textbf{prepare}: Verify role prerequisites within the container
    
    \item \textbf{converge}: Run ansible-playbook against the container -- this is where the role's tasks are executed
    
    \item \textbf{test/verify}: Optional pytest sequence runs inside the container to validate convergence
\end{enumerate}

\subsection{Example: DNS Role Testing}

The DNS role molecule test confirmed:
\begin{itemize}
    \item Podman creates an Ubuntu 20.04 container (hostname "server1")
    
    \item Bind9 packages install correctly
    
    \item /etc/bind directory is created with named.conf.options written
    
    \item Molecule platform name was changed from "molecule-ubuntu" to "server1" to match the role's hostname gate
    
    \item Group assignment updated to "dns\_servers" for proper role targeting
\end{itemize}

\subsection{Known Limitations}

Some roles cannot be fully tested on molecule containers:

\begin{itemize}
    \item \textbf{nfs}: Hostname gates (\texttt{'snowy' in inventory\_hostname}) and hardcoded mount paths depend on real NFS servers
    
    \item \textbf{openbsd}: The role targets OpenBSD hosts specifically; molecule tests run on Ubuntu containers by default
    
    \item \textbf{jetson-nano}: Hardware-specific (NVIDIA Jetson); molecule tests only verify file placement, not actual provisioning
\end{itemize}

\subsection{Running Tests Manually}

To test a specific role from the stargate host:
\begin{lstlisting}[style=mystyle]
cd /mnt/clusterfs2/workspace/lab-franklin/ansible/collections/ansible_collections/lab/franklin/roles/<role>/molecule
molecule converge --scenario-name default
    \end{lstlisting}

To verify convergence results:
\begin{lstlisting}[style=mystyle]
molecule test --scenario-name default
    \end{lstlisting}

If the podman driver is unavailable, fall back to docker:
\begin{lstlisting}[style=mystyle]
molecule converge --driver-name docker
    \end{lstlisting}

\subsection{Test Results Dashboard}

Testing output logs are captured in the molecule scenario directory under \texttt{logs/}. The full lab-franklin test suite (all roles with molecule tests) is run via the top-level Makefile.am \texttt{test} target:
\begin{lstlisting}[style=mystyle]
make test
# Runs ansible-lint --all-roles and molecule converge per role
    \end{lstlisting}

\subsection{pytest-ansible Transition Plan}

As molecule is being superseded, new tests should use pytest-ansible where possible:
\begin{itemize}
    \item Migrate the test directory structure from \texttt{molecule/} to a \texttt{tests/} directory with pytest configuration
    
    \item Use ansible-test for collection sanity tests (built into ansible-core)
    
    \item Keep molecule only where full environment isolation is needed (e.g., DNS zone file validation, NTP time sync testing)
\end{itemize}

\subsection{Pending Tasks and Todos}

\begin{itemize}
    \item \textbf{TODO: Re-establish internal PyPI server on stargate}. The pypi\_internal role is ready (role exists, molecule test added Aug 2026) but the actual PyPI service needs to be provisioned on the storage host. Requires: NFS mount of /mnt/storage1 available on stargate, web server (nginx/apache) installed and configured, pip packages populated into /mnt/storage1/LAB/pypi with PEP 503 structure. Run: \texttt{ansible-playbook ... -l storage --tags pypi\_internal} once the infrastructure is in place.
\end{itemize}

\section{Ansible Automation}

\subsection{Current Landscape (August 2026)}

As of August 2026, the Ansible ecosystem is split between two distribution paths:

\begin{itemize}
    \item \textbf{Ansible community package} -- the traditional monolithic install (pip install ansible). Current version is 13.x (built on ansible-core 2.20). This includes 85+ collections with thousands of modules and plugins bundled together. Follows semantic versioning, releases two major versions per year plus monthly minor releases.
    
    \item \textbf{ansible-core} -- the minimal core engine. Maintains three versions at a time (current + two older). Versioned independently from the community package (2.19, 2.20, 2.21). Releases every four weeks for minor updates.

\end{itemize}

\textbf{EOL versions as of August 2026}: Ansible 10--12 and ansible-core 2.15--2.18 are EOL. Do not use these in new work. Only Ansible 13.x / ansible-core 2.20 is current; 14.0.0/2.21 is in development.

\subsection{Molecule Testing -- Status and Deprecation}

Molecule (the canonical integration testing framework for Ansible roles and playbooks) remains \textbf{active} as of August 2026 but its role has shifted:

\begin{itemize}
    \item Molecule v26.4.x is the latest series, released June 2026. It still supports only the latest two major versions of Ansible (N/N-1).
    
    \item Molecule is being superseded as the preferred test strategy by:
    \begin{itemize}
        \item \texttt{pytest-ansible} -- pytest extension for testing Ansible module/plugin Python code directly.
        
        \item \texttt{tox-ansible} -- integration with tox to run tests across multiple Python interpreters and ansible-core versions simultaneously.
        
        \item Execution Environments (containers) as the primary test target, tested via \texttt{ansible-navigator}.
    \end{itemize}
    
    \item The \textbf{community.molecule} collection still exists but new projects should prefer pytest-ansible + tox-ansible for role/test harness work. Molecule is not yet officially deprecated -- it is maintained on a best-effort basis while the community transitions.
    
    \item For existing molecule test harnesses (such as those in lab-franklin's Ansible collection at \texttt{ansible/collections/ansible\_collections/lab/franklin/roles/}), they will continue to work with ansible-core 2.20 and pytest-6.x but should be audited for:
    \begin{itemize}
        \item Podman driver compatibility (molecule continues using podman or docker as default driver)
        \item Python version requirements (molecule requires Python $\ge$3.9)
        \item Plugin deprecations in newer collections
    \end{itemize}

\end{itemize}

\textbf{Note}: For lab-franklin's molecule tests, keep testing with podman for now. The container-based approach is still valid; only the \textit{preferred tooling around it} is shifting toward pytest-ansible + execution environments.

\chapter{lab/franklin Collection}

The lab/franklin Ansible collection lives at:
\begin{lstlisting}[style=mystyle]
/mnt/clusterfs2/workspace/lab-franklin/ansible/collections/ansible_collections/lab/franklin/
    \end{lstlisting}

It is the central automation artifact for the entire lab.bitsmasher.net infrastructure, managing provisioning, configuration, and testing across all hosts.

\section{Collection Manifest (galaxy.yml)}

The collection follows standard Ansible Galaxy conventions with namespace \texttt{lab} and name \texttt{franklin}:

\begin{itemize}
    \item \textbf{Version}: 1.0.0
    
    \item \textbf{License}: MIT
    
    \item \textbf{Author}: franklin <franklin@bitsmasher.net>
    
    \item \textbf{Repository}: \url{https://github.com/devsecfranklin/lab-franklin}
    
    \item \textbf{Documentation}: \url{https://github.com/devsecfranklin/lab-franklin/docs}
\end{itemize}

\section{Role Inventory}

The collection contains roles organized under the roles/ directory. As of August 2026, the following roles are maintained with molecule test harnesses:

\begin{lstlisting}[style=mystyle]
roles/
├── apt_mirror/          # APT mirror configuration
├── beagleboard/         # BeagleBoard device provisioning
├── common/              # Shared base configuration (bootstrap.sh, common utils)
├── container_registry/  # Docker/container registry setup
├── ctfd/                # CTFd platform deployment
├── desktop/             # Desktop environment provisioning
├── docker/              # Docker engine installation
├── documentation/       # Documentation generation roles
├── dhcp/                # DHCP server configuration (dhcpd.conf)
├── dns/                 # BIND/named DNS server
├── golang/              # Go language toolchain
├── jetson-nano/         # NVIDIA Jetson Nano provisioning
├── k3s_agent/           # k3s worker node
├── k3s_server/          # k3s control plane server
├── k8s/                 # Kubernetes generic setup
├── latex/               # LaTeX environment
├── logging/             # Centralized logging
├── music/               # Music service configuration
├── nfs/                 # NFS server/client
├── nix/                 # Nix package manager
├── openbsd/             # OpenBSD host bootstrap (blowfish target)
│   ├── tasks/main.yml   # Includes files.yml, timezone setup
│   └── files/           # bootstrap.sh, login.conf, hosts entries
├── paloalto/            # Palo Alto firewall config
├── prereq/              # Prerequisite packages
├── pypi_internal/       # Internal PyPI repository setup
│   ├── tasks/main.yml   # Symlink pypi_dir -> html_dir/pypi
│   └── vars/main.yml    # html_dir=/var/www/html, pypi_dir=/mnt/storage1/LAB/pypi
├── python/              # Python environment
├── raspberrypi/         # Raspberry Pi provisioning
├── samba/               # Samba file sharing
├── security/            # Hardening and security baselines
├── shell/               # Shell configuration
├── ssh/                 # SSH server hardening
├── tls/                 # TLS certificate management
├── chonk/               # chonk-specific configuration
├── cluster/             # Cluster-wide settings
    \end{lstlisting}

Roles currently \textbf{missing molecule tests}: edge, mail, odroid, website (incomplete roles). These have tasks/main.yml but no test harness.

\section{Workspace Structure}

The full project layout:
\begin{lstlisting}[style=mystyle]
/workspace/lab-franklin/
├── Makefile.am                  # Top-level autotools target (Python, Docker, dev)
├── configure.ac                 # Autotools config (autoconf)
├── bootstrap.sh                 # Cross-platform bootstrapper
├── network_update.sh            # One-shot maintenance script (hardened 2026-08-02)
├── ansible/                     # ANSIBLE_HOME root
│   ├── playbook.yml             # Main playbook
│   ├── hosts                    # Inventory file
│   └── collections/             # Collections path
│       ├── ansible_collections/ # Galaxy namespace
│           └── lab/franklin/    # The lab/franklin collection
│               ├── galaxy.yml
│               ├── roles/       # All role modules (37 total)
│               └── docs/latex/  # LaTeX docs build (autotools)
├── container/                   # Docker/Podman configs
├── terraform/                   # Terraform infrastructure definitions
├── bin/                         # Shared utility scripts (common.sh)
└── docs/manual/                 # Lab manual documentation (this document)
    \end{lstlisting}

\section{Build System (Autotools)}

The project uses GNU Autotools for the root-level build:

\begin{itemize}
    \item \textbf{configure.ac}: Defines package (lab-franklin v0.2), libtool support, Python $\ge$ 3.9 requirement, podman/ansible/gpg checks, git version tagging
    
    \item \textbf{Makefile.am}: Top-level targets include test (ansible-lint + molecule), security (venv bootstrap), dev (Python venv in BUILDDIR)
    
    \item \textbf{bootstrap.sh}: Cross-platform bootstrapper that detects OS (Debian/RedHat/OpenBSD/macOS/Linux) and installs platform-specific packages, then runs aclocal $\rightarrow$ autoreconf $\rightarrow$ automake $\rightarrow$ configure
    
    \item \textbf{docs/manual/Makefile.am}: Standalone autotools stub for LaTeX docs build (clean target removes *.aux files)
\end{itemize}

The \texttt{network\_update.sh} script was hardened on 2026-08-02 with: set -euo pipefail, auto-resolved paths, pinned K3s version validation, ANSIBLE_HOME environment checks, removed dangerous clush/apt-get blast radius, and dead code removal. It runs the main playbook at \texttt{ansible/playbook.yml} as its primary action.

\chapter{Testing from Stargate}

\section{Molecule Test Harnesses on stargate.lab.bitsmasher.net}

The stargate host (10.10.16.66, Debian 12 bookworm) serves as the primary testing platform for the lab/franklin Ansible collection. Testing is run from \texttt{/mnt/clusterfs2/workspace/lab-franklin/ansible/collections/ansible\_collections/lab/franklin/roles/\textless role\textgreater /molecule/}.

\subsection{Test Environment}

Molecule tests on stargate use:
\begin{itemize}
    \item \textbf{Driver}: podman (containers, not Docker -- avoids daemon dependency)
    
    \item \textbf{Base image}: Ubuntu 20.04 for most roles, with role-specific variations
    
    \item \textbf{Execution user}: openclaw user (has sudo -n, no password required)
    
    \item \textbf{Python}: $\ge$ 3.9 (required by both ansible and molecule)
\end{itemize}

\subsection{Test Execution Flow}

For a given role, the molecule test cycle is:
\begin{enumerate}
    \item \textbf{dependency}: Install collection dependencies from galaxy.yml
    
    \item \textbf{syntax}: Validate playbook YAML syntax
    
    \item \textbf{create}: Launch podman container (e.g., "server1") as the test target host
    
    \item \textbf{prepare}: Verify role prerequisites within the container
    
    \item \textbf{converge}: Run ansible-playbook against the container -- this is where the role's tasks are executed
    
    \item \textbf{test/verify}: Optional pytest sequence runs inside the container to validate convergence
\end{enumerate}

\subsection{Example: DNS Role Testing}

The DNS role molecule test confirmed:
\begin{itemize}
    \item Podman creates an Ubuntu 20.04 container (hostname "server1")
    
    \item Bind9 packages install correctly
    
    \item /etc/bind directory is created with named.conf.options written
    
    \item Molecule platform name was changed from "molecule-ubuntu" to "server1" to match the role's hostname gate
    
    \item Group assignment updated to "dns\_servers" for proper role targeting
\end{itemize}

\subsection{Known Limitations}

Some roles cannot be fully tested on molecule containers:

\begin{itemize}
    \item \textbf{nfs}: Hostname gates (\texttt{'snowy' in inventory\_hostname}) and hardcoded mount paths depend on real NFS servers
    
    \item \textbf{openbsd}: The role targets OpenBSD hosts specifically; molecule tests run on Ubuntu containers by default
    
    \item \textbf{jetson-nano}: Hardware-specific (NVIDIA Jetson); molecule tests only verify file placement, not actual provisioning
\end{itemize}

\subsection{Running Tests Manually}

To test a specific role from the stargate host:
\begin{lstlisting}[style=mystyle]
cd /mnt/clusterfs2/workspace/lab-franklin/ansible/collections/ansible_collections/lab/franklin/roles/<role>/molecule
molecule converge --scenario-name default
    \end{lstlisting}

To verify convergence results:
\begin{lstlisting}[style=mystyle]
molecule test --scenario-name default
    \end{lstlisting}

If the podman driver is unavailable, fall back to docker:
\begin{lstlisting}[style=mystyle]
molecule converge --driver-name docker
    \end{lstlisting}

\subsection{Test Results Dashboard}

Testing output logs are captured in the molecule scenario directory under \texttt{logs/}. The full lab-franklin test suite (all roles with molecule tests) is run via the top-level Makefile.am \texttt{test} target:
\begin{lstlisting}[style=mystyle]
make test
# Runs ansible-lint --all-roles and molecule converge per role
    \end{lstlisting}

\subsection{pytest-ansible Transition Plan}

As molecule is being superseded, new tests should use pytest-ansible where possible:
\begin{itemize}
    \item Migrate the test directory structure from \texttt{molecule/} to a \texttt{tests/} directory with pytest configuration
    
    \item Use ansible-test for collection sanity tests (built into ansible-core)
    
    \item Keep molecule only where full environment isolation is needed (e.g., DNS zone file validation, NTP time sync testing)
\end{itemize}

\subsection{Pending Tasks and Todos}

\begin{itemize}
    \item \textbf{TODO: Re-establish internal PyPI server on stargate}. The pypi\_internal role is ready (role exists, molecule test added Aug 2026) but the actual PyPI service needs to be provisioned on the storage host. Requires: NFS mount of /mnt/storage1 available on stargate, web server (nginx/apache) installed and configured, pip packages populated into /mnt/storage1/LAB/pypi with PEP 503 structure. Run: \texttt{ansible-playbook ... -l storage --tags pypi\_internal} once the infrastructure is in place.
\end{itemize}


\chapter{Additional Infrastructure}

\section{Blowfish -- Database Host}
blowfish is the planned database host (OpenBSD). Currently unreachable; SSH key not yet authorized.

% \section{Blowfish -- Database Host}

\subsection{Host Identity}

\begin{itemize}
    \item \textbf{Hostname}: blowfish.lab.bitsmasher.net (alias: blowfish)
    
    \item \textbf{IP}: 10.10.12.15
    
    \item \textbf{MAC}: 98:b7:85:21:ad:77 (from DHCP reservation)
    
    \item \textbf{OS}: OpenBSD (planned; current DNS/DHCP entry references it as an OpenBSD host)
\end{itemize}

\subsection{Role in the Lab}

Blowfish is provisioned as a \textbf{database host} within the lab infrastructure. It was intended to serve as a dedicated database server, separate from the primary infrastructure hosts (chonk, stargate, skynet).

The OpenBSD role in the lab-franklin Ansible collection includes configuration references specific to blowfish:
\begin{itemize}
    \item Login configuration uses \texttt{blowfish} password hashing algorithm (\texttt{localcipher=blowfish,a}) in \texttt{/etc/login.conf}.
    
    \item The bootstrap.sh for OpenBSD hosts includes a direct \texttt{/etc/hosts} entry for blowfish (10.10.12.15).
    
    \item DHCP reservation is configured in the lab's dhcpd.conf role.
\end{itemize}

\subsection{Current Status}

Blowfish has been \textbf{unreachable} from chonk since the initial infrastructure audit. The IP 10.10.12.15 resolved from DNS (stargate's /etc/hosts), but authentication was denied -- no matching SSH key was authorized on that host.

Key references:
\begin{itemize}
    \item network\_update.sh: blowfish is the default target for the \texttt{openbsd\_bootstrap()} function.
    
    \item OpenBSD role tasks: includes commented-out ansible setup command for blowfish (indicating the host was expected to be online).
\end{itemize}

The host may need network investigation -- it resolved from DNS, had a DHCP reservation, and MAC address on file, suggesting it was once part of the lab but has since gone offline or been decommissioned.

\subsection{Planned Configuration (OpenBSD)}

When brought online, the expected setup for blowfish is:

\begin{itemize}
    \item SSH key-based authentication via \texttt{~/.ssh/authorized\_keys}
    
    \item OpenBSD's built-in \texttt{pf} firewall with custom rules
    
    \item Database software stack (likely PostgreSQL or MariaDB on OpenBSD)
    
    \item Ansible-managed configuration via the lab-franklin collection, using the openbsd role as the base profile.
\end{itemize}

\section{Blowfish -- Database Host}

\subsection{Host Identity}

\begin{itemize}
    \item \textbf{Hostname}: blowfish.lab.bitsmasher.net (alias: blowfish)
    
    \item \textbf{IP}: 10.10.12.15
    
    \item \textbf{MAC}: 98:b7:85:21:ad:77 (from DHCP reservation)
    
    \item \textbf{OS}: OpenBSD (planned; current DNS/DHCP entry references it as an OpenBSD host)
\end{itemize}

\subsection{Role in the Lab}

Blowfish is provisioned as a \textbf{database host} within the lab infrastructure. It was intended to serve as a dedicated database server, separate from the primary infrastructure hosts (chonk, stargate, skynet).

The OpenBSD role in the lab-franklin Ansible collection includes configuration references specific to blowfish:
\begin{itemize}
    \item Login configuration uses \texttt{blowfish} password hashing algorithm (\texttt{localcipher=blowfish,a}) in \texttt{/etc/login.conf}.
    
    \item The bootstrap.sh for OpenBSD hosts includes a direct \texttt{/etc/hosts} entry for blowfish (10.10.12.15).
    
    \item DHCP reservation is configured in the lab's dhcpd.conf role.
\end{itemize}

\subsection{Current Status}

Blowfish has been \textbf{unreachable} from chonk since the initial infrastructure audit. The IP 10.10.12.15 resolved from DNS (stargate's /etc/hosts), but authentication was denied -- no matching SSH key was authorized on that host.

Key references:
\begin{itemize}
    \item network\_update.sh: blowfish is the default target for the \texttt{openbsd\_bootstrap()} function.
    
    \item OpenBSD role tasks: includes commented-out ansible setup command for blowfish (indicating the host was expected to be online).
\end{itemize}

The host may need network investigation -- it resolved from DNS, had a DHCP reservation, and MAC address on file, suggesting it was once part of the lab but has since gone offline or been decommissioned.

\subsection{Planned Configuration (OpenBSD)}

When brought online, the expected setup for blowfish is:

\begin{itemize}
    \item SSH key-based authentication via \texttt{~/.ssh/authorized\_keys}
    
    \item OpenBSD's built-in \texttt{pf} firewall with custom rules
    
    \item Database software stack (likely PostgreSQL or MariaDB on OpenBSD)
    
    \item Ansible-managed configuration via the lab-franklin collection, using the openbsd role as the base profile.
\end{itemize}


\section{Minecraft Server}
The Minecraft server runs on skynet, with world data stored on NFS and version-controlled via bare git repos.

% \section{Minecraft Server}

\subsection{Host and Location}

The Minecraft server (Dreamland) runs on \textbf{chonk.lab.bitsmasher.net} (10.10.8.60), a primary compute host with RX 9070 / Navi 48 GPU. Server files are at /home/minecraft.

Previous host was skynet.lab.bitsmasher.net (10.10.16.10) during the CraftBukkit/early Paper era.

\subsection{Current Configuration}

The server runs \textbf{Paper 1.21.11} on Java 21 (Debian Bookworm). Key settings:
\begin{itemize}
    \item Gamemode: Survival
    \item Difficulty: Normal
    \item Max players: 69
    \item PvP: Disabled
    \item View distance: 10 (simulation distance: 12)
    \item Server port: 25565
\end{itemize}

The server name is \textit{Dreamland} (MOTD: \texttt{\&aDreamland\&7 | \&fEst. 2012}).

\subsection{Core Plugins}

\begin{itemize}
    \item \textbf{EssentialsX} -- Core, Chat, AntiBuild
    \item \textbf{LuckPerms} -- Permissions (backed by MariaDB)
    \item \textbf{CoreProtect} -- Audit/log (backed by MySQL)
\end{itemize}

\subsection{Interactive Map}

BlueMap provides an interactive 3D web map: \url{https://chonk.lab.bitsmasher.net:8100}

\subsection{Storage}

World data is approximately 13 GB. BlueMap assets are ~6.5 GB. All files on local disk (no NFS dependency for world data).

\textbf{NO GIT-LFS}: Chunk files and world data must not use git-lfs. Use bare git repos for any versioned world metadata.

\subsection{Administrative Access}

Operations automation via start.sh with a custom rcon-based controller. Management infrastructure ports: server 25565, BlueMap 8100, Plan analytics 8804.

\textbf{Ops team}: theDevilsVoice (level 4), Perryman (level 4).

\textbf{Whitelist}: slyborg4realz, UnvaluedShoe79, BeefyGamer76y, some_garlic, KindredFiori, secove_, DrewPellino, NotPebl3, VibyCat, CallMeTesla.

\section{Minecraft Server}

\subsection{Host and Location}

The Minecraft server (Dreamland) runs on \textbf{chonk.lab.bitsmasher.net} (10.10.8.60), a primary compute host with RX 9070 / Navi 48 GPU. Server files are at /home/minecraft.

Previous host was skynet.lab.bitsmasher.net (10.10.16.10) during the CraftBukkit/early Paper era.

\subsection{Current Configuration}

The server runs \textbf{Paper 1.21.11} on Java 21 (Debian Bookworm). Key settings:
\begin{itemize}
    \item Gamemode: Survival
    \item Difficulty: Normal
    \item Max players: 69
    \item PvP: Disabled
    \item View distance: 10 (simulation distance: 12)
    \item Server port: 25565
\end{itemize}

The server name is \textit{Dreamland} (MOTD: \texttt{\&aDreamland\&7 | \&fEst. 2012}).

\subsection{Core Plugins}

\begin{itemize}
    \item \textbf{EssentialsX} -- Core, Chat, AntiBuild
    \item \textbf{LuckPerms} -- Permissions (backed by MariaDB)
    \item \textbf{CoreProtect} -- Audit/log (backed by MySQL)
\end{itemize}

\subsection{Interactive Map}

BlueMap provides an interactive 3D web map: \url{https://chonk.lab.bitsmasher.net:8100}

\subsection{Storage}

World data is approximately 13 GB. BlueMap assets are ~6.5 GB. All files on local disk (no NFS dependency for world data).

\textbf{NO GIT-LFS}: Chunk files and world data must not use git-lfs. Use bare git repos for any versioned world metadata.

\subsection{Administrative Access}

Operations automation via start.sh with a custom rcon-based controller. Management infrastructure ports: server 25565, BlueMap 8100, Plan analytics 8804.

\textbf{Ops team}: theDevilsVoice (level 4), Perryman (level 4).

\textbf{Whitelist}: slyborg4realz, UnvaluedShoe79, BeefyGamer76y, some_garlic, KindredFiori, secove_, DrewPellino, NotPebl3, VibyCat, CallMeTesla.


\chapter{Disaster Recovery}

% \section{Disaster Recovery}

\subsection{Critical Dependencies Chain}

The lab infrastructure has a cascading dependency chain. A failure in one layer propagates:

\begin{enumerate}
    \item \textbf{Time (NTP)} $\rightarrow$ if time drift exceeds 128 seconds, Kerberos tickets expire and TLS handshakes fail
    
    \item \textbf{DNS (ns1/BIND)} $\rightarrow$ without DNS resolution, all hostname-based services (Kerberos SRV records, NTP upstream, LDAP) break
    
    \item \textbf{Kerberos (odroid-c1/KDC)} $\rightarrow$ without the KDC, authentication fails for all Kerberos-authenticated services
    
    \item \textbf{LDAP (bbb1/slapd)} $\rightarrow$ without LDAP, directory lookups fail and user/service account information is unavailable
\end{enumerate}

Recovery priority should follow this chain in reverse: fix time first, then DNS, then KDC, then LDAP.

\subsection{Failure Scenarios and Recovery}

\subsubsection{NTP Failure -- All Hosts Lose Time Sync}

Symptoms: \texttt{timedatectl} shows "System clock synchronized: no", Kerberos tickets rejected.

Recovery:
\begin{enumerate}
    \item Verify GPS unit is locked on odroid-c1 (time host)
    
    \item Check ntpd/ntpsec service: \texttt{systemctl status ntpsec}
    
    \item If GPS is offline, enable holdover mode (internal oscillator continues at reduced accuracy)
    
    \item On affected clients, force resync: \texttt{ntpdate -s time.lab.bitsmasher.net} or wait for ntpd's step-sync to kick in
    
    \item Verify with \texttt{ntpq -p} on clients and \texttt{ntptime} on the server
\end{enumerate}

\subsubsection{DNS (ns1) Offline -- Cascade Failure}

Symptoms: hostnames don't resolve, Kerberos SRV discovery fails.

Recovery:
\begin{enumerate}
    \item Check ns1 reachability: ping from any known-good host
    
    \item If ns1 is up but not answering: check BIND (named) service status
    
    \item If BIND has crashed due to config error: restore from last known-good named.conf
    
    \item As a workaround, add manual entries to /etc/hosts on affected hosts until DNS recovers
    
    \item Update krb5.conf with explicit \texttt{kdc = odroid-c1.lab.bitsmasher.net}
\end{enumerate}

\subsubsection{LDAP (bbb1) Down -- No Directory Service}

Symptoms: slapd failed, status.sh reports both slapd and ldapsearch as failing.

Recovery:
\begin{enumerate}
    \item SSH to bbb1/lab.bitsmasher.net as root
    
    \item Check TLS certificate validity in the configured cert path
    
    \item Regenerate or replace expired certificates
    
    \item Restart slapd: \texttt{systemctl restart slapd}
    
    \item Verify with ldapsearch for franklin and sly DNs
\end{enumerate}

\subsubsection{KDC (odroid-c1) Offline -- No Authentication}

Symptoms: kinit fails for all realms, service authentication denied.

Recovery:
\begin{enumerate}
    \item SSH to odroid-c1 as root (franklin user's key is not authorized)
    
    \item Check kdc process: \texttt{kadmind/krb5kdc status}
    
    \item Verify KDC database integrity: \texttt{kdb5_util list}
    
    \item Restart kerberos services if needed
    
    \item Regenerate host keytabs for affected services via \texttt{kadmin.local}
\end{enumerate}

\subsection{Backup and Restoration}

\begin{itemize}
    \item Ansible playbooks (lab-franklin collection) are the single source of truth -- they can reconstruct any configured state
    
    \item Bare git repos with GPG-encrypted pass entries handle credential backup
    
    \item KDC database (\texttt{\$KRB5\_KDB\_FILE}) should be backed up regularly on odroid-c1
    
    \item DNS zone files should have local copies (in the DNS role's files/ directory)
\end{itemize}

\subsection{Prevention: Monitoring}

The current monitoring approach is manual via status.sh. For improved resilience, consider:
\begin{itemize}
    \item Automated ping/SSH checks on critical hosts in HEARTBEAT.md
    
    \item Nagios/Zabbix-style monitoring on a dedicated host
    
    \item Email/slack alerts for service down conditions
\end{itemize}

The lab-franklin Ansible collection should eventually include a monitoring role to automate this.

\section{Disaster Recovery}

\subsection{Critical Dependencies Chain}

The lab infrastructure has a cascading dependency chain. A failure in one layer propagates:

\begin{enumerate}
    \item \textbf{Time (NTP)} $\rightarrow$ if time drift exceeds 128 seconds, Kerberos tickets expire and TLS handshakes fail
    
    \item \textbf{DNS (ns1/BIND)} $\rightarrow$ without DNS resolution, all hostname-based services (Kerberos SRV records, NTP upstream, LDAP) break
    
    \item \textbf{Kerberos (odroid-c1/KDC)} $\rightarrow$ without the KDC, authentication fails for all Kerberos-authenticated services
    
    \item \textbf{LDAP (bbb1/slapd)} $\rightarrow$ without LDAP, directory lookups fail and user/service account information is unavailable
\end{enumerate}

Recovery priority should follow this chain in reverse: fix time first, then DNS, then KDC, then LDAP.

\subsection{Failure Scenarios and Recovery}

\subsubsection{NTP Failure -- All Hosts Lose Time Sync}

Symptoms: \texttt{timedatectl} shows "System clock synchronized: no", Kerberos tickets rejected.

Recovery:
\begin{enumerate}
    \item Verify GPS unit is locked on odroid-c1 (time host)
    
    \item Check ntpd/ntpsec service: \texttt{systemctl status ntpsec}
    
    \item If GPS is offline, enable holdover mode (internal oscillator continues at reduced accuracy)
    
    \item On affected clients, force resync: \texttt{ntpdate -s time.lab.bitsmasher.net} or wait for ntpd's step-sync to kick in
    
    \item Verify with \texttt{ntpq -p} on clients and \texttt{ntptime} on the server
\end{enumerate}

\subsubsection{DNS (ns1) Offline -- Cascade Failure}

Symptoms: hostnames don't resolve, Kerberos SRV discovery fails.

Recovery:
\begin{enumerate}
    \item Check ns1 reachability: ping from any known-good host
    
    \item If ns1 is up but not answering: check BIND (named) service status
    
    \item If BIND has crashed due to config error: restore from last known-good named.conf
    
    \item As a workaround, add manual entries to /etc/hosts on affected hosts until DNS recovers
    
    \item Update krb5.conf with explicit \texttt{kdc = odroid-c1.lab.bitsmasher.net}
\end{enumerate}

\subsubsection{LDAP (bbb1) Down -- No Directory Service}

Symptoms: slapd failed, status.sh reports both slapd and ldapsearch as failing.

Recovery:
\begin{enumerate}
    \item SSH to bbb1/lab.bitsmasher.net as root
    
    \item Check TLS certificate validity in the configured cert path
    
    \item Regenerate or replace expired certificates
    
    \item Restart slapd: \texttt{systemctl restart slapd}
    
    \item Verify with ldapsearch for franklin and sly DNs
\end{enumerate}

\subsubsection{KDC (odroid-c1) Offline -- No Authentication}

Symptoms: kinit fails for all realms, service authentication denied.

Recovery:
\begin{enumerate}
    \item SSH to odroid-c1 as root (franklin user's key is not authorized)
    
    \item Check kdc process: \texttt{kadmind/krb5kdc status}
    
    \item Verify KDC database integrity: \texttt{kdb5_util list}
    
    \item Restart kerberos services if needed
    
    \item Regenerate host keytabs for affected services via \texttt{kadmin.local}
\end{enumerate}

\subsection{Backup and Restoration}

\begin{itemize}
    \item Ansible playbooks (lab-franklin collection) are the single source of truth -- they can reconstruct any configured state
    
    \item Bare git repos with GPG-encrypted pass entries handle credential backup
    
    \item KDC database (\texttt{\$KRB5\_KDB\_FILE}) should be backed up regularly on odroid-c1
    
    \item DNS zone files should have local copies (in the DNS role's files/ directory)
\end{itemize}

\subsection{Prevention: Monitoring}

The current monitoring approach is manual via status.sh. For improved resilience, consider:
\begin{itemize}
    \item Automated ping/SSH checks on critical hosts in HEARTBEAT.md
    
    \item Nagios/Zabbix-style monitoring on a dedicated host
    
    \item Email/slack alerts for service down conditions
\end{itemize}

The lab-franklin Ansible collection should eventually include a monitoring role to automate this.


\chapter{Troubleshooting Guide}

% \section{Troubleshooting Guide}

\subsection{SSH Connection Failures}

\subsubsection{Host Key Mismatch}

Symptom: \texttt{@@@@@ WARNING: REMOTE HOST IDENTIFICATION HAS CHANGED! @@@@@}

Fix:
\begin{lstlisting}[style=mystyle]
ssh-keygen -R <hostname>
    \end{lstlisting}

This occurs when a host was re-imaged with the same IP. Check \texttt{/etc/ssh/ssh\_config.d/} for stale entries.

\subsubsection{Auth Denied (Port 22 Open, Key Rejected)}

Symptom: TCP port 22 is reachable but SSH returns "Permission denied".

Common causes:
\begin{itemize}
    \item No matching key in the remote host's authorized\_keys
    
    \item The key file permissions are too open on the target (must be mode 600 for .pub and 700 for ~/.ssh)
    
    \item Wrong user context (e.g., trying franklin@ when only root is configured)
\end{itemize}

Fix: Add the openclaw ed25519 key to the target host's authorized\_keys file. The pubkey is:
\begin{lstlisting}[style=mystyle]
ssh-ed25519 AAAAC3NzaC1lZDI1NTE5AAAAIJb9mV02PpxD8VpzYCnu7192dTwMnGWUc3qh55BIeElN openclaw@chonk
    \end{lstlisting}

\subsubsection{No Route to Host}

Symptom: SSH hangs or returns "No route to host". The host is likely powered off. Common in Jetson nodes (node902, etc.) that are manually managed.

\subsubsection{GSSAPI Interference}

Symptom: SSH times out waiting for GSSAPI authentication before falling back to key auth.

Fix: Add to \texttt{~/.ssh/config}:
\begin{lstlisting}[style=mystyle]
Host *
    GSSAPIAuthentication no
    IdentitiesOnly yes
    IdentityFile ~/.ssh/id_ed25519_openclaw
    \end{lstlisting}

\subsection{Service-Specific Diagnostics}

\subsubsection{Kerberos Authentication Failure}

\begin{tabularx}{\textwidth}{l X}
\textbf{Symptom} & \textbf{Likely Cause} \\
\hline
kinit fails with "clock skew" & NTP not synced -- check time.host first\\
kinit fails with "KDC has no support for encryption type" & Wrong krb5.conf realm or KDC unavailable\\
kadmin.local works but kadmin (remote) fails & Port 749 firewall rule on odroid-c1\\
Ticket expired after reboot & Time drift during host downtime
\end{tabularx}

Verify: \texttt{klist}, \texttt{kinit username@LAB.BITSMASHER.NET}, \texttt{ktutil} for keytabs.

\subsubsection{NTP Issues}

\begin{tabularx}{\textwidth}{l X}
\textbf{Symptom} & \textbf{Likely Cause} \\
\hline
ntpd not running & ntpsec service not enabled\\
"no server suitable for synchronization" & No upstream reachable (time host offline)\\
Clock gradually drifting & Missing driftfile or Jetson oscillator drift\\
Time jumps by several seconds & Large correction applied; step-sync mode engaged
\end{tabularx}

Verify: \texttt{ntpq -p}, \texttt{timedatectl}, check /var/log/ntpstats/.

\subsubsection{NFS Mount Issues}

\begin{tabularx}{\textwidth}{l X}
\textbf{Symptom} & \textbf{Likely Cause} \\
\hline
mount hangs & NFS server unreachable or export doesn't list client\\
"Stale file handle" & Server-side data was deleted/recreated\\
Permission denied on mount & /etc/exports restricts the client IP
\end{tabularx}

Verify: \texttt{showmount -e <server>}, check NFS server's exports, verify client IP matches the export rule.

\subsection{Common Ansible Troubleshooting}

\subsubsection{Module Not Found / Missing Collection}

\begin{lstlisting}[style=mystyle]
ansible-galaxy collection install lab.franklin
# or for individual modules:
ansible-galaxy collection install community.general
    \end{lstlisting}

\subsubsection{Molecule Test Failures on New Hosts}

If molecule tests fail on a new host, the most common causes are:
\begin{enumerate}
    \item Podman not installed or permissions wrong (openclaw user needs podman group)
    
    \item Docker driver configured instead of podman -- edit \texttt{molecule.yml} to set driver
    
    \item Python version mismatch (molecule requires $\ge$3.9)
    
    \item Role hostname gates preventing test execution (e.g., the DNS role's hostname check for the real server vs. molecule container name)
\end{enumerate}

\subsection{General Debugging Checklist}

When a service is down, work through this checklist:

\begin{enumerate}
    \item Can you ping the host? ($\rightarrow$ network up?)
    
    \item Is SSH responding? ($\rightarrow$ host OS running?)
    
    \item Does the service process exist? (\texttt{systemctl status or ps aux})
    
    \item Are logs showing errors? (\texttt{journalctl -u <service> --since today})
    
    \item Can you connect to the service port from another lab host? (\texttt{nc -zv <host> <port>})
    
    \item Is DNS resolving correctly? (\texttt{dig or getent hosts})
    
    \item Is time synchronized? (\texttt{timedatectl})
\end{enumerate}

\subsection{Quick Reference Commands}

\begin{lstlisting}[style=mystyle]
# Check if a host is reachable from chonk
ping -c 1 <hostname>
ssh -o ConnectTimeout=5 openclaw@<host> "echo alive"

# Check service status remotely
ssh openclaw@<host> "systemctl status <service>"

# Verify time sync across hosts
for host in stargate skynet chonk; do
    ssh openclaw@$host "timedatectl | grep synchronized"
done

# Verify DNS resolution from a client
dig @time.lab.bitsmasher.net <hostname>

# Check Kerberos auth
kinit username@LAB.BITSMASHER.NET
klist
    \end{lstlisting}

\section{Troubleshooting Guide}

\subsection{SSH Connection Failures}

\subsubsection{Host Key Mismatch}

Symptom: \texttt{@@@@@ WARNING: REMOTE HOST IDENTIFICATION HAS CHANGED! @@@@@}

Fix:
\begin{lstlisting}[style=mystyle]
ssh-keygen -R <hostname>
    \end{lstlisting}

This occurs when a host was re-imaged with the same IP. Check \texttt{/etc/ssh/ssh\_config.d/} for stale entries.

\subsubsection{Auth Denied (Port 22 Open, Key Rejected)}

Symptom: TCP port 22 is reachable but SSH returns "Permission denied".

Common causes:
\begin{itemize}
    \item No matching key in the remote host's authorized\_keys
    
    \item The key file permissions are too open on the target (must be mode 600 for .pub and 700 for ~/.ssh)
    
    \item Wrong user context (e.g., trying franklin@ when only root is configured)
\end{itemize}

Fix: Add the openclaw ed25519 key to the target host's authorized\_keys file. The pubkey is:
\begin{lstlisting}[style=mystyle]
ssh-ed25519 AAAAC3NzaC1lZDI1NTE5AAAAIJb9mV02PpxD8VpzYCnu7192dTwMnGWUc3qh55BIeElN openclaw@chonk
    \end{lstlisting}

\subsubsection{No Route to Host}

Symptom: SSH hangs or returns "No route to host". The host is likely powered off. Common in Jetson nodes (node902, etc.) that are manually managed.

\subsubsection{GSSAPI Interference}

Symptom: SSH times out waiting for GSSAPI authentication before falling back to key auth.

Fix: Add to \texttt{~/.ssh/config}:
\begin{lstlisting}[style=mystyle]
Host *
    GSSAPIAuthentication no
    IdentitiesOnly yes
    IdentityFile ~/.ssh/id_ed25519_openclaw
    \end{lstlisting}

\subsection{Service-Specific Diagnostics}

\subsubsection{Kerberos Authentication Failure}

\begin{tabularx}{\textwidth}{l X}
\textbf{Symptom} & \textbf{Likely Cause} \\
\hline
kinit fails with "clock skew" & NTP not synced -- check time.host first\\
kinit fails with "KDC has no support for encryption type" & Wrong krb5.conf realm or KDC unavailable\\
kadmin.local works but kadmin (remote) fails & Port 749 firewall rule on odroid-c1\\
Ticket expired after reboot & Time drift during host downtime
\end{tabularx}

Verify: \texttt{klist}, \texttt{kinit username@LAB.BITSMASHER.NET}, \texttt{ktutil} for keytabs.

\subsubsection{NTP Issues}

\begin{tabularx}{\textwidth}{l X}
\textbf{Symptom} & \textbf{Likely Cause} \\
\hline
ntpd not running & ntpsec service not enabled\\
"no server suitable for synchronization" & No upstream reachable (time host offline)\\
Clock gradually drifting & Missing driftfile or Jetson oscillator drift\\
Time jumps by several seconds & Large correction applied; step-sync mode engaged
\end{tabularx}

Verify: \texttt{ntpq -p}, \texttt{timedatectl}, check /var/log/ntpstats/.

\subsubsection{NFS Mount Issues}

\begin{tabularx}{\textwidth}{l X}
\textbf{Symptom} & \textbf{Likely Cause} \\
\hline
mount hangs & NFS server unreachable or export doesn't list client\\
"Stale file handle" & Server-side data was deleted/recreated\\
Permission denied on mount & /etc/exports restricts the client IP
\end{tabularx}

Verify: \texttt{showmount -e <server>}, check NFS server's exports, verify client IP matches the export rule.

\subsection{Common Ansible Troubleshooting}

\subsubsection{Module Not Found / Missing Collection}

\begin{lstlisting}[style=mystyle]
ansible-galaxy collection install lab.franklin
# or for individual modules:
ansible-galaxy collection install community.general
    \end{lstlisting}

\subsubsection{Molecule Test Failures on New Hosts}

If molecule tests fail on a new host, the most common causes are:
\begin{enumerate}
    \item Podman not installed or permissions wrong (openclaw user needs podman group)
    
    \item Docker driver configured instead of podman -- edit \texttt{molecule.yml} to set driver
    
    \item Python version mismatch (molecule requires $\ge$3.9)
    
    \item Role hostname gates preventing test execution (e.g., the DNS role's hostname check for the real server vs. molecule container name)
\end{enumerate}

\subsection{General Debugging Checklist}

When a service is down, work through this checklist:

\begin{enumerate}
    \item Can you ping the host? ($\rightarrow$ network up?)
    
    \item Is SSH responding? ($\rightarrow$ host OS running?)
    
    \item Does the service process exist? (\texttt{systemctl status or ps aux})
    
    \item Are logs showing errors? (\texttt{journalctl -u <service> --since today})
    
    \item Can you connect to the service port from another lab host? (\texttt{nc -zv <host> <port>})
    
    \item Is DNS resolving correctly? (\texttt{dig or getent hosts})
    
    \item Is time synchronized? (\texttt{timedatectl})
\end{enumerate}

\subsection{Quick Reference Commands}

\begin{lstlisting}[style=mystyle]
# Check if a host is reachable from chonk
ping -c 1 <hostname>
ssh -o ConnectTimeout=5 openclaw@<host> "echo alive"

# Check service status remotely
ssh openclaw@<host> "systemctl status <service>"

# Verify time sync across hosts
for host in stargate skynet chonk; do
    ssh openclaw@$host "timedatectl | grep synchronized"
done

# Verify DNS resolution from a client
dig @time.lab.bitsmasher.net <hostname>

# Check Kerberos auth
kinit username@LAB.BITSMASHER.NET
klist
    \end{lstlisting}


\backmatter

\chapter{Infrastructure Services}
\section{Network Time (NTP)}

\subsection{Overview}

Time synchronization is one of the most critical and least-visible aspects of network management. All lab services -- Kerberos authentication, TLS certificate validation, log correlation, distributed databases, and filesystem consistency -- depend on time accuracy. The lab uses a two-tier NTP architecture:

\begin{enumerate}
    \item \textbf{Stratum 1 source}: GPS-referenced clock on the \texttt{time} host (10.10.12.2) provides the lab's authoritative time reference.
    
    \item \textbf{Stratum 2+ clients}: All other hosts sync to \texttt{time.lab.bitsmasher.net}.
\end{enumerate}

\subsection{The Time Host}

\begin{itemize}
    \item \textbf{Hostname}: time.lab.bitsmasher.net (alias: time)
    
    \item \textbf{IP}: 10.10.12.2
    
    \item \textbf{User}: openclaw
    
    \item \textbf{OS}: Debian 12 bookworm
    
    \item \textbf{Role}: Stratum-1 NTP server / reference clock source for the lab
\end{itemize}

The time host has a GPS unit connected (serial/tty interface) providing UTC discipline. Its ntpd is configured via ntpsec to serve the lab subnet and use the GPS PPS signal as the primary time reference.

\subsection{Configuration on Client Hosts}

The NTP configuration is managed by Ansible via the \texttt{ntp} role in the lab-franklin collection. Here is what a typical client configuration looks like (as seen on stargate):

\begin{lstlisting}[style=mystyle]
# /etc/ntpsec/ntp.conf -- managed by Ansible
driftfile /var/lib/ntp/ntp.drift
leapfile /usr/share/zoneinfo/leap-seconds.list

statsdir /var/log/ntpstats/

statistics loopstats peerstats clockstats
filegen loopstats file loopstats type day enable
filegen peerstats file peerstats type day enable
filegen clockstats file clockstats type day enable

server time.lab.bitsmasher.net
restrict 10.10.8.0/21
restrict -4 default
\end{lstlisting}

Key points:
\begin{itemize}
    \item \texttt{driftfile}: Tracks oscillator drift between NTP restarts. Critical for stability on hosts that reboot frequently (Jetson devices, test VMs).
    
    \item \texttt{statistics}: Logs loop/peer/clock stats daily in /var/log/ntpstats/. Use these to verify sync health.
    
    \item \texttt{server}: Points to the lab's time host -- never pool.ntp.org for infrastructure hosts that need deterministic time sources.
    
    \item \texttt{restrict}: 
    \begin{itemize}
        \item \texttt{10.10.8.0/21} grants query access to the main lab subnet (used by other internal clients or monitoring tools).
        
        \item \texttt{-4 default} blocks all IPv4 queries not explicitly allowed, providing a deny-by-default posture.
    \end{itemize}
\end{itemize}

\subsection{Installation and Setup Steps}

For a new host that needs NTP:

\begin{enumerate}
    \item Install ntpsec packages:
    \begin{lstlisting}[style=mystyle]
apt update && apt install -y ntpsec
    \end{lstlisting}

    \item Configure /etc/ntpsec/ntp.conf (either manually or via Ansible ntp role):
    \begin{enumerate}
        \item Set \texttt{server time.lab.bitsmasher.net} as the upstream reference.
        
        \item Set timezone: \texttt{timedatectl set-timezone America/Denver} (or appropriate TZ).
        
        \item Create the log directory if it doesn't exist: \texttt{mkdir -p /var/log/ntpstats}.
    \end{enumerate}

    \item Start and enable the service:
    \begin{lstlisting}[style=mystyle]
systemctl enable --now ntpsec.service
timedatectl  # verify "System clock synchronized: yes"
    \end{lstlisting}
\end{enumerate}

\subsection{Verification on Clients}

Check synchronization status with these commands:

\begin{lstlisting}[style=mystyle]
# Check if the system clock is synchronized
timedatectl | grep "System clock synchronized"

# Check NTP peer status (if ntpq is available)
ntpq -p

# View recent sync logs
journalctl -u ntpsec --since today
    \end{lstlisting}

\subsection{Notes for Deployment}

\begin{itemize}
    \item The Ansible \texttt{ntp} role handles all of the above automatically across lab hosts. It installs packages, copies configuration, sets timezone, and creates log directories.
    
    \item For Jetson devices (node900--node903), drift compensation is especially important -- they boot frequently and have less stable oscillators than desktop/server hardware.
    
    \item The GPS unit on the time host should be verified periodically: confirm it's locking to satellites, not using its internal oscillator as a holdover source. Check this by running \texttt{ntptime} on the time host itself.
    
    \item If the time host goes offline and all clients lose sync beyond their maximum correction threshold (typically 128 seconds), systems may need manual time correction before Kerberos/TLS will function again.
\end{itemize}

\section{Storage Infrastructure}

\subsection{NFS Mounts}

The lab relies on NFS for shared storage across hosts. Key mount points:

\begin{itemize}
    \item \textbf{/mnt/clusterfs2}: Primary cluster workspace mounted on chonk and stargate (from the NFS server). This is where Ansible roles, molecule tests, and project files live.
    
    \item \textbf{/mnt/snowy}: Formerly a full NFS share; now just a mounted drive after the hard disk was physically removed. Should be audited for stale mount entries.
\end{itemize}

The NFS role in lab-franklin's Ansible collection handles:
\begin{itemize}
    \item Server-side export configuration (/etc/exports)
    \item Client-side fstab management and mount points
    \item Mount path conventions per role (previously /mnt/clusterfs, /mnt/backup1, /mnt/storage{1,2,3})
\end{itemize}

Note: The NFS role has hostname gates (\texttt{'snowy' in inventory\_hostname}) that prevent testing on molecule containers. Hardcoded mount paths depend on real NFS servers being online. This is a known limitation of the test harness.

\subsection{Snapshot and Backup Strategy}

The lab uses bare git repos for version control and backup:
\begin{itemize}
    \item GPG-encrypted repositories synced across hosts for credential protection (Stash House project)
    \item Pass-based password store with GPG backend
    \item No LFS -- ABSOLUTE RULE: do not use git-lfs anywhere. Microsoft once held Minecraft chunk files hostage when a repo hit the 10GB GitHub limit with LFS enabled.
\end{itemize}

\subsection{Local Storage Considerations}

\begin{itemize}
    \item Jetson nodes (node90x): limited eMMC storage; rely on NFS for workspace data
    \item chonk: large local disk, primary build and gateway host
    \item stargate: Debian 12 box with full Ansible workspace mirror
\end{itemize}

\subsection{LFS Policy Reminder}

\textbf{ABSOLUTE RULE: Do NOT use git-lfs anywhere in any lab-franklin repository.}

This applies across all repos: training-ai, minecraft, Ansible collections, everything. The incident with Minecraft chunk files being held hostage by GitHub when the repo hit 10GB with LFS is a real precedent -- never risk it.


\chapter{File Services: NFS}
\section{What We Know About NFS}
\section{Storage Infrastructure}

\subsection{NFS Mounts}

The lab relies on NFS for shared storage across hosts. Key mount points:

\begin{itemize}
    \item \textbf{/mnt/clusterfs2}: Primary cluster workspace mounted on chonk and stargate (from the NFS server). This is where Ansible roles, molecule tests, and project files live.
    
    \item \textbf{/mnt/snowy}: Formerly a full NFS share; now just a mounted drive after the hard disk was physically removed. Should be audited for stale mount entries.
\end{itemize}

The NFS role in lab-franklin's Ansible collection handles:
\begin{itemize}
    \item Server-side export configuration (/etc/exports)
    \item Client-side fstab management and mount points
    \item Mount path conventions per role (previously /mnt/clusterfs, /mnt/backup1, /mnt/storage{1,2,3})
\end{itemize}

Note: The NFS role has hostname gates (\texttt{'snowy' in inventory\_hostname}) that prevent testing on molecule containers. Hardcoded mount paths depend on real NFS servers being online. This is a known limitation of the test harness.

\subsection{Snapshot and Backup Strategy}

The lab uses bare git repos for version control and backup:
\begin{itemize}
    \item GPG-encrypted repositories synced across hosts for credential protection (Stash House project)
    \item Pass-based password store with GPG backend
    \item No LFS -- ABSOLUTE RULE: do not use git-lfs anywhere. Microsoft once held Minecraft chunk files hostage when a repo hit the 10GB GitHub limit with LFS enabled.
\end{itemize}

\subsection{Local Storage Considerations}

\begin{itemize}
    \item Jetson nodes (node90x): limited eMMC storage; rely on NFS for workspace data
    \item chonk: large local disk, primary build and gateway host
    \item stargate: Debian 12 box with full Ansible workspace mirror
\end{itemize}

\subsection{LFS Policy Reminder}

\textbf{ABSOLUTE RULE: Do NOT use git-lfs anywhere in any lab-franklin repository.}

This applies across all repos: training-ai, minecraft, Ansible collections, everything. The incident with Minecraft chunk files being held hostage by GitHub when the repo hit 10GB with LFS is a real precedent -- never risk it.


\section{How to Test the NFS Ansible Role}
% Testing the NFS Ansible Role

\section{Overview}

The NFS role at \texttt{ansible/collections/ansible\_collections/lab/franklin/roles/nfs} has a molecule test harness, but testing is complex due to hostname-dependent task gates. This chapter documents how to run tests and what to expect.

\section{Prerequisites}

\begin{itemize}
  \item podman available on the test host (chonk tested successfully)
  \item The lab.franklin collection with all dependencies installed
  \item A molecule environment with an Ubuntu-based container (the role expects Debian/Ubuntu apt packages)
\end{itemize}

\section{Running the Test}

From the role directory:

\begin{verbatim}
cd /mnt/clusterfs2/workspace/lab-franklin/ansible/collections/ansible_collections/lab/franklin/roles/nfs
molecule converge --scenario-name default
molecule verify --scenario-name default
\end{verbatim}

\section{What the Test Validates}

\begin{itemize}
  \item \textbf{Package installation}: \texttt{nfs-kernel-server}, \texttt{nfs-common}, \texttt{rpcbind}, \texttt{nfs-client} are installed
  \item \textbf{Configuration file deployment}: \texttt{/etc/default/nfs-kernel-server}, \texttt{/etc/default/nfs-common}, \texttt{/etc/idmapd.conf} are written with correct content
  \item \textbf{Directory creation}: Export mount points (\texttt{/mnt/clusterfs}, \texttt{/mnt/storage\{1,2,3\}}) are created
  \item \textbf{Export template rendering}: Jinja2 templates generate valid \texttt{/etc/exports} content
\end{itemize}

\section{Current Test Limitations}

The role has two categories of gates that block full testing:

\subsection*{1. Server-side hostname gates (in tasks/main.yml)}

\begin{verbatim}
when: 'snowy' in inventory_hostname
when: 'thelio' in inventory_hostname
\end{verbatim}

Molecule containers cannot satisfy these conditions, so \texttt{nfs\_server\_*.yml} task files never execute. The server-side converge phase is effectively skipped.

\subsection*{2. molecule environment variable guards}

Most mount and service-start tasks check:
\begin{verbatim}
when: HOMELAB_MOLECULE_TEST is not defined
\end{verbatim}

This prevents actual mounts from happening during testing (which would require real NFS servers), but it also means the \texttt{nfs\_client.yml} tasks that matter most for client-side validation never run their key stages.

\section{How to Fix the Test Harness}

To make molecule tests meaningful, the following changes are needed:

\begin{enumerate}
  \item \textbf{Add test-specific inventory}: In \texttt{molecule/default/inventory.yml}, define a hostgroup (e.g., \texttt{nfs\_servers}) that includes container hostnames matching the role's hostname gates. Example:
  \begin{verbatim}
    all:
      children:
        nfs_servers:
          hosts:
            server1:
              ansible_connection: docker
        nfs_clients:
          hosts:
            client1:
              ansible_connection: docker
  \end{verbatim}

  \item \textbf{Add override variables in molecule\_prepare.yml}: Set \texttt{nfs\_exports\_snowy} and \texttt{nfs\_exports\_thelio} to container-safe values that use local paths, so export template rendering can be validated without real storage.

  \item \textbf{Create a test scenario for client-only testing}: A second molecule scenario with only nfs\_clients inventory that tests mount directory creation, fstab line insertion, and idmapd.conf deployment. Use \texttt{state: present} instead of \texttt{state: mounted} in the molecule overrides so it doesn't try to connect to a live NFS server.

  \item \textbf{Replace hostname gates with role variables}: Replace \texttt{'snowy' in inventory\_hostname} with a variable like \texttt{nfs\_server\_mode | default(false)} so tests can enable server tasks without depending on real hostnames.

  \item \textbf{Add idempotence tests}: A second converge pass to verify that running the role again produces no changes (the \texttt{force: no} on idmapd.conf copy already helps here).
\end{enumerate}

\section{Integration Test Targets}

An integration test target exists at \texttt{tests/integration/targets/nfs-server/} but it is minimal. For thorough testing, add:

\begin{itemize}
  \item Export verification (check \texttt{showmount -e})
  \item Mount/unmount cycle tests on client container
  \item Kerberos security validation (try mounting with and without \texttt{sec=krb5i})
  \item Permission checks after mount (verify root\_squash behavior)
\end{itemize}

\section{Molecule Test Checklist for Future Work}

Before declaring the NFS role test-ready:

\begin{itemize}
  \item [ ] Server-side converge passes (with hostname override in molecule inventory)
  \item [ ] Client-side converge passes (mount directory creation, fstab entries)
  \item [ ] Verify phase checks exports file content and configuration files
  \item [ ] Idempotence test: second converge shows zero changed tasks
  \item [ ] Integration test validates live mount on container
\end{itemize}



\chapter{Automation \& Tooling}
\section{Ansible Automation}

\subsection{Current Landscape (August 2026)}

As of August 2026, the Ansible ecosystem is split between two distribution paths:

\begin{itemize}
    \item \textbf{Ansible community package} -- the traditional monolithic install (pip install ansible). Current version is 13.x (built on ansible-core 2.20). This includes 85+ collections with thousands of modules and plugins bundled together. Follows semantic versioning, releases two major versions per year plus monthly minor releases.
    
    \item \textbf{ansible-core} -- the minimal core engine. Maintains three versions at a time (current + two older). Versioned independently from the community package (2.19, 2.20, 2.21). Releases every four weeks for minor updates.

\end{itemize}

\textbf{EOL versions as of August 2026}: Ansible 10--12 and ansible-core 2.15--2.18 are EOL. Do not use these in new work. Only Ansible 13.x / ansible-core 2.20 is current; 14.0.0/2.21 is in development.

\subsection{Molecule Testing -- Status and Deprecation}

Molecule (the canonical integration testing framework for Ansible roles and playbooks) remains \textbf{active} as of August 2026 but its role has shifted:

\begin{itemize}
    \item Molecule v26.4.x is the latest series, released June 2026. It still supports only the latest two major versions of Ansible (N/N-1).
    
    \item Molecule is being superseded as the preferred test strategy by:
    \begin{itemize}
        \item \texttt{pytest-ansible} -- pytest extension for testing Ansible module/plugin Python code directly.
        
        \item \texttt{tox-ansible} -- integration with tox to run tests across multiple Python interpreters and ansible-core versions simultaneously.
        
        \item Execution Environments (containers) as the primary test target, tested via \texttt{ansible-navigator}.
    \end{itemize}
    
    \item The \textbf{community.molecule} collection still exists but new projects should prefer pytest-ansible + tox-ansible for role/test harness work. Molecule is not yet officially deprecated -- it is maintained on a best-effort basis while the community transitions.
    
    \item For existing molecule test harnesses (such as those in lab-franklin's Ansible collection at \texttt{ansible/collections/ansible\_collections/lab/franklin/roles/}), they will continue to work with ansible-core 2.20 and pytest-6.x but should be audited for:
    \begin{itemize}
        \item Podman driver compatibility (molecule continues using podman or docker as default driver)
        \item Python version requirements (molecule requires Python $\ge$3.9)
        \item Plugin deprecations in newer collections
    \end{itemize}

\end{itemize}

\textbf{Note}: For lab-franklin's molecule tests, keep testing with podman for now. The container-based approach is still valid; only the \textit{preferred tooling around it} is shifting toward pytest-ansible + execution environments.

\chapter{lab/franklin Collection}

The lab/franklin Ansible collection lives at:
\begin{lstlisting}[style=mystyle]
/mnt/clusterfs2/workspace/lab-franklin/ansible/collections/ansible_collections/lab/franklin/
    \end{lstlisting}

It is the central automation artifact for the entire lab.bitsmasher.net infrastructure, managing provisioning, configuration, and testing across all hosts.

\section{Collection Manifest (galaxy.yml)}

The collection follows standard Ansible Galaxy conventions with namespace \texttt{lab} and name \texttt{franklin}:

\begin{itemize}
    \item \textbf{Version}: 1.0.0
    
    \item \textbf{License}: MIT
    
    \item \textbf{Author}: franklin <franklin@bitsmasher.net>
    
    \item \textbf{Repository}: \url{https://github.com/devsecfranklin/lab-franklin}
    
    \item \textbf{Documentation}: \url{https://github.com/devsecfranklin/lab-franklin/docs}
\end{itemize}

\section{Role Inventory}

The collection contains roles organized under the roles/ directory. As of August 2026, the following roles are maintained with molecule test harnesses:

\begin{lstlisting}[style=mystyle]
roles/
├── apt_mirror/          # APT mirror configuration
├── beagleboard/         # BeagleBoard device provisioning
├── common/              # Shared base configuration (bootstrap.sh, common utils)
├── container_registry/  # Docker/container registry setup
├── ctfd/                # CTFd platform deployment
├── desktop/             # Desktop environment provisioning
├── docker/              # Docker engine installation
├── documentation/       # Documentation generation roles
├── dhcp/                # DHCP server configuration (dhcpd.conf)
├── dns/                 # BIND/named DNS server
├── golang/              # Go language toolchain
├── jetson-nano/         # NVIDIA Jetson Nano provisioning
├── k3s_agent/           # k3s worker node
├── k3s_server/          # k3s control plane server
├── k8s/                 # Kubernetes generic setup
├── latex/               # LaTeX environment
├── logging/             # Centralized logging
├── music/               # Music service configuration
├── nfs/                 # NFS server/client
├── nix/                 # Nix package manager
├── openbsd/             # OpenBSD host bootstrap (blowfish target)
│   ├── tasks/main.yml   # Includes files.yml, timezone setup
│   └── files/           # bootstrap.sh, login.conf, hosts entries
├── paloalto/            # Palo Alto firewall config
├── prereq/              # Prerequisite packages
├── pypi_internal/       # Internal PyPI repository setup
│   ├── tasks/main.yml   # Symlink pypi_dir -> html_dir/pypi
│   └── vars/main.yml    # html_dir=/var/www/html, pypi_dir=/mnt/storage1/LAB/pypi
├── python/              # Python environment
├── raspberrypi/         # Raspberry Pi provisioning
├── samba/               # Samba file sharing
├── security/            # Hardening and security baselines
├── shell/               # Shell configuration
├── ssh/                 # SSH server hardening
├── tls/                 # TLS certificate management
├── chonk/               # chonk-specific configuration
├── cluster/             # Cluster-wide settings
    \end{lstlisting}

Roles currently \textbf{missing molecule tests}: edge, mail, odroid, website (incomplete roles). These have tasks/main.yml but no test harness.

\section{Workspace Structure}

The full project layout:
\begin{lstlisting}[style=mystyle]
/workspace/lab-franklin/
├── Makefile.am                  # Top-level autotools target (Python, Docker, dev)
├── configure.ac                 # Autotools config (autoconf)
├── bootstrap.sh                 # Cross-platform bootstrapper
├── network_update.sh            # One-shot maintenance script (hardened 2026-08-02)
├── ansible/                     # ANSIBLE_HOME root
│   ├── playbook.yml             # Main playbook
│   ├── hosts                    # Inventory file
│   └── collections/             # Collections path
│       ├── ansible_collections/ # Galaxy namespace
│           └── lab/franklin/    # The lab/franklin collection
│               ├── galaxy.yml
│               ├── roles/       # All role modules (37 total)
│               └── docs/latex/  # LaTeX docs build (autotools)
├── container/                   # Docker/Podman configs
├── terraform/                   # Terraform infrastructure definitions
├── bin/                         # Shared utility scripts (common.sh)
└── docs/manual/                 # Lab manual documentation (this document)
    \end{lstlisting}

\section{Build System (Autotools)}

The project uses GNU Autotools for the root-level build:

\begin{itemize}
    \item \textbf{configure.ac}: Defines package (lab-franklin v0.2), libtool support, Python $\ge$ 3.9 requirement, podman/ansible/gpg checks, git version tagging
    
    \item \textbf{Makefile.am}: Top-level targets include test (ansible-lint + molecule), security (venv bootstrap), dev (Python venv in BUILDDIR)
    
    \item \textbf{bootstrap.sh}: Cross-platform bootstrapper that detects OS (Debian/RedHat/OpenBSD/macOS/Linux) and installs platform-specific packages, then runs aclocal $\rightarrow$ autoreconf $\rightarrow$ automake $\rightarrow$ configure
    
    \item \textbf{docs/manual/Makefile.am}: Standalone autotools stub for LaTeX docs build (clean target removes *.aux files)
\end{itemize}

The \texttt{network\_update.sh} script was hardened on 2026-08-02 with: set -euo pipefail, auto-resolved paths, pinned K3s version validation, ANSIBLE_HOME environment checks, removed dangerous clush/apt-get blast radius, and dead code removal. It runs the main playbook at \texttt{ansible/playbook.yml} as its primary action.

\chapter{Testing from Stargate}

\section{Molecule Test Harnesses on stargate.lab.bitsmasher.net}

The stargate host (10.10.16.66, Debian 12 bookworm) serves as the primary testing platform for the lab/franklin Ansible collection. Testing is run from \texttt{/mnt/clusterfs2/workspace/lab-franklin/ansible/collections/ansible\_collections/lab/franklin/roles/\textless role\textgreater /molecule/}.

\subsection{Test Environment}

Molecule tests on stargate use:
\begin{itemize}
    \item \textbf{Driver}: podman (containers, not Docker -- avoids daemon dependency)
    
    \item \textbf{Base image}: Ubuntu 20.04 for most roles, with role-specific variations
    
    \item \textbf{Execution user}: openclaw user (has sudo -n, no password required)
    
    \item \textbf{Python}: $\ge$ 3.9 (required by both ansible and molecule)
\end{itemize}

\subsection{Test Execution Flow}

For a given role, the molecule test cycle is:
\begin{enumerate}
    \item \textbf{dependency}: Install collection dependencies from galaxy.yml
    
    \item \textbf{syntax}: Validate playbook YAML syntax
    
    \item \textbf{create}: Launch podman container (e.g., "server1") as the test target host
    
    \item \textbf{prepare}: Verify role prerequisites within the container
    
    \item \textbf{converge}: Run ansible-playbook against the container -- this is where the role's tasks are executed
    
    \item \textbf{test/verify}: Optional pytest sequence runs inside the container to validate convergence
\end{enumerate}

\subsection{Example: DNS Role Testing}

The DNS role molecule test confirmed:
\begin{itemize}
    \item Podman creates an Ubuntu 20.04 container (hostname "server1")
    
    \item Bind9 packages install correctly
    
    \item /etc/bind directory is created with named.conf.options written
    
    \item Molecule platform name was changed from "molecule-ubuntu" to "server1" to match the role's hostname gate
    
    \item Group assignment updated to "dns\_servers" for proper role targeting
\end{itemize}

\subsection{Known Limitations}

Some roles cannot be fully tested on molecule containers:

\begin{itemize}
    \item \textbf{nfs}: Hostname gates (\texttt{'snowy' in inventory\_hostname}) and hardcoded mount paths depend on real NFS servers
    
    \item \textbf{openbsd}: The role targets OpenBSD hosts specifically; molecule tests run on Ubuntu containers by default
    
    \item \textbf{jetson-nano}: Hardware-specific (NVIDIA Jetson); molecule tests only verify file placement, not actual provisioning
\end{itemize}

\subsection{Running Tests Manually}

To test a specific role from the stargate host:
\begin{lstlisting}[style=mystyle]
cd /mnt/clusterfs2/workspace/lab-franklin/ansible/collections/ansible_collections/lab/franklin/roles/<role>/molecule
molecule converge --scenario-name default
    \end{lstlisting}

To verify convergence results:
\begin{lstlisting}[style=mystyle]
molecule test --scenario-name default
    \end{lstlisting}

If the podman driver is unavailable, fall back to docker:
\begin{lstlisting}[style=mystyle]
molecule converge --driver-name docker
    \end{lstlisting}

\subsection{Test Results Dashboard}

Testing output logs are captured in the molecule scenario directory under \texttt{logs/}. The full lab-franklin test suite (all roles with molecule tests) is run via the top-level Makefile.am \texttt{test} target:
\begin{lstlisting}[style=mystyle]
make test
# Runs ansible-lint --all-roles and molecule converge per role
    \end{lstlisting}

\subsection{pytest-ansible Transition Plan}

As molecule is being superseded, new tests should use pytest-ansible where possible:
\begin{itemize}
    \item Migrate the test directory structure from \texttt{molecule/} to a \texttt{tests/} directory with pytest configuration
    
    \item Use ansible-test for collection sanity tests (built into ansible-core)
    
    \item Keep molecule only where full environment isolation is needed (e.g., DNS zone file validation, NTP time sync testing)
\end{itemize}

\subsection{Pending Tasks and Todos}

\begin{itemize}
    \item \textbf{TODO: Re-establish internal PyPI server on stargate}. The pypi\_internal role is ready (role exists, molecule test added Aug 2026) but the actual PyPI service needs to be provisioned on the storage host. Requires: NFS mount of /mnt/storage1 available on stargate, web server (nginx/apache) installed and configured, pip packages populated into /mnt/storage1/LAB/pypi with PEP 503 structure. Run: \texttt{ansible-playbook ... -l storage --tags pypi\_internal} once the infrastructure is in place.
\end{itemize}

\section{Ansible-Linting Infrastructure}

All ansible roles in the lab collection should be linted on a regular basis to maintain configuration quality and catch issues early. The lab uses three tools: \texttt{ansible-lint}, \texttt{yamllint}, and \texttt{ansible-test sanity}.

\subsection{Tools and Configuration}

Each linter is configured in the \texttt{test/} directory:
\begin{itemize}
    \item \texttt{.ansible-lint.yml} — ansible-lint configuration with warnings for deprecated modules and experimental rules disabled
    \item \texttt{.yamllint} / \texttt{.yamllintrc} — YAML linting extends default rules with lab-specific indentation and spacing constraints
\end{itemize}

\subsection{Coverage Status (44 Roles)}

Based on the most recent audit of \texttt{ansible/collections/ansible\_collections/lab/franklin/roles/}:

\begin{description}
    \item[Complete (30 roles)] apt-mirror, beagleboard, chonk, cluster, common, container-registry, ctfd, desktop, dhcp, dns, docker, documentation, golang, k3s-agent, k3s-server, kerberos, ldap, logging, minecraft, music, nfs, nix, ntp, openbsd, paloalto, prereq, pypi-internal, python, raspberrypi, ssh
    \item[Minimal/Functional (9 roles)] media, samba, shell, tls, website, apt-mirror-stub, jetson-nano, k8s (partial), dns-stub
    \item[Stubs (4 roles)] edge, extensions, latex, odroid
\end{description}

Roles with the smallest role definitions tend to have the fewest linting issues. The most problematic roles are usually those that include dynamic template generation or custom module invocation.

\subsection{Proposed Linting Pipeline}

A new \texttt{test/lint\_all.sh} script should be created to run all three linters in sequence:

\begin{verbatim}
# Example lint_all.sh structure:
#!/bin/bash
set -euo pipefail

cd "$(dirname "$0")/.."

echo "=== ansible-lint ==="
ansible-lint roles/ 2>&1 || echo "ansible-lint warnings found"

echo "=== yamllint ==="
yamllint -c test/.yamllint roles/ 2>&1 || echo "yamllint issues found"

echo "=== ansible-test sanity ==="
ansible-test sanity --docker default -v 2>&1 || echo "sanity check failures"
\end{verbatim}

The linting pipeline should be integrated into:
\begin{itemize}
    \item \textbf{GitHub Actions:} Run on PR and nightly builds (similar to existing bandit/trivy/tfsec workflows)
    \item \textbf{Hourly cron job on stargate:} Optional for CI-style validation without external triggers
\end{itemize}

\subsection{Linting Exceptions by Role}

Some roles require special handling during linting:
\begin{itemize}
    \item \texttt{tls}: Includes shell commands with \texttt{kss -kfe} which may trigger warning 403 (deprecated module) — explicitly skip in config
    \item \texttt{samba}: Uses custom \texttt{smb.conf.j2} templates that require validation via \texttt{testparm} post-converge
    \item \texttt{k8s}: Contains static manifests rather than role tasks — should be moved to a shared files directory or documented as a reference role
\end{itemize}

% Chapter on Terraform Infrastructure Management
% Add this as % Chapter on Terraform Infrastructure Management
% Add this as % Chapter on Terraform Infrastructure Management
% Add this as \input{sections/terraform-infra} in lab-manual.tex after the Ansible chapter

\chapter{Terraform Infrastructure Management}

Terraform configurations in this lab manage cloud infrastructure across two providers: DigitalOcean and Google Cloud Platform. All configurations are version-controlled and include automated security scanning via GitHub Actions.

\section{DigitalOcean Infrastructure}

The \texttt{terraform/digital\_ocean/} directory manages the external-facing web presence for bitsmasher.net.

\subsection{Provisioned Resources}

\begin{itemize}
    \item \textbf{Droplet:} 512MB Debian 12 in lon1 region (www.bitsmasher.net)
    \item \textbf{Domain:} bitsmasher.net with full DNS management
    \item \textbf{DNS Records:} A, MX, TXT (SPF/DMARC/DKIM), NS, and ACME challenge records
\end{itemize}

\subsection{Key DNS Records}

\begin{center}
\begin{tabular}{lll}
\hline
Name & Type & Purpose \\ \hline
www.bitsmasher.net & A & Web server host (178.62.60.55) \\
@ & MX & mail.protonmail.ch 10 (ProtonMail routing) \\
@ & TXT & SPF and DMARC policies \\
protonmail.\_domainkey & CNAME & DKIM email verification \\
\_acme-challenge.*.lab & TXT & Let's Encrypt certificate challenges \\ \hline
\end{tabular}
\end{center}

\subsection{Configuration Notes}

The droplet includes:
\begin{itemize}
    \item Automated backups enabled (immutable against accidental deletion)
    \item SSH key injection via terraform data source
    \item No external monitoring (consider enabling DO agent for uptime tracking)
\end{itemize}

Provider authentication uses a DigitalOcean API token managed through pass/direnv integration. The token is passed at runtime, never committed to the repository.

\section{Google Cloud Platform Infrastructure}

The \texttt{terraform/google/} directory manages GCP resources for Stash House free-tier operations.

\subsection{Resources Managed}

\begin{itemize}
    \item \textbf{VPC Network:} \texttt{stash-house-vpc} dedicated network
    \item \textbf{Compute Instance:} e2-micro Debian 12 (Always Free Tier eligible) in us-central1-a
    \item \textbf{Firewall Rule:} \texttt{allow-ssh} opening port 22 to the VPC
    \item \textbf{Billing Budget Alert:} Tripwire alert at \$0.01 threshold for free-tier monitoring
\end{itemize}

The GCP configuration is intentionally minimal, designed around Google's Always Free Tier limits to maintain zero-cost infrastructure.

\section{Security Scanning}

All Terraform configurations are automatically scanned by two GitHub Actions workflows:
\begin{itemize}
    \item \texttt{tfsec-sarif.yml} — IaC security scanning (generates SARIF reports)
    \item \texttt{tfsec-comment.yml} — Posts security findings as PR comments
\end{itemize}

Additionally, the DigitalOcean directory includes existing \texttt{README.md} with doctl and pass integration instructions for manual operations.

\section{Testing}

Terraform configurations are validated through:
\begin{itemize}
    \item Go-based integration tests (\texttt{terraform/test/main\_test.go})
    \item Plan dry-run validation via GitHub Actions (markdown.yml workflow)
    \item Manual state reconciliation checks (compare terraform state against live infrastructure)
\end{itemize}

State drift detection is recommended as a regular maintenance task — compare \texttt{terraform state list} output against actual cloud resource inventories quarterly.
 in lab-manual.tex after the Ansible chapter

\chapter{Terraform Infrastructure Management}

Terraform configurations in this lab manage cloud infrastructure across two providers: DigitalOcean and Google Cloud Platform. All configurations are version-controlled and include automated security scanning via GitHub Actions.

\section{DigitalOcean Infrastructure}

The \texttt{terraform/digital\_ocean/} directory manages the external-facing web presence for bitsmasher.net.

\subsection{Provisioned Resources}

\begin{itemize}
    \item \textbf{Droplet:} 512MB Debian 12 in lon1 region (www.bitsmasher.net)
    \item \textbf{Domain:} bitsmasher.net with full DNS management
    \item \textbf{DNS Records:} A, MX, TXT (SPF/DMARC/DKIM), NS, and ACME challenge records
\end{itemize}

\subsection{Key DNS Records}

\begin{center}
\begin{tabular}{lll}
\hline
Name & Type & Purpose \\ \hline
www.bitsmasher.net & A & Web server host (178.62.60.55) \\
@ & MX & mail.protonmail.ch 10 (ProtonMail routing) \\
@ & TXT & SPF and DMARC policies \\
protonmail.\_domainkey & CNAME & DKIM email verification \\
\_acme-challenge.*.lab & TXT & Let's Encrypt certificate challenges \\ \hline
\end{tabular}
\end{center}

\subsection{Configuration Notes}

The droplet includes:
\begin{itemize}
    \item Automated backups enabled (immutable against accidental deletion)
    \item SSH key injection via terraform data source
    \item No external monitoring (consider enabling DO agent for uptime tracking)
\end{itemize}

Provider authentication uses a DigitalOcean API token managed through pass/direnv integration. The token is passed at runtime, never committed to the repository.

\section{Google Cloud Platform Infrastructure}

The \texttt{terraform/google/} directory manages GCP resources for Stash House free-tier operations.

\subsection{Resources Managed}

\begin{itemize}
    \item \textbf{VPC Network:} \texttt{stash-house-vpc} dedicated network
    \item \textbf{Compute Instance:} e2-micro Debian 12 (Always Free Tier eligible) in us-central1-a
    \item \textbf{Firewall Rule:} \texttt{allow-ssh} opening port 22 to the VPC
    \item \textbf{Billing Budget Alert:} Tripwire alert at \$0.01 threshold for free-tier monitoring
\end{itemize}

The GCP configuration is intentionally minimal, designed around Google's Always Free Tier limits to maintain zero-cost infrastructure.

\section{Security Scanning}

All Terraform configurations are automatically scanned by two GitHub Actions workflows:
\begin{itemize}
    \item \texttt{tfsec-sarif.yml} — IaC security scanning (generates SARIF reports)
    \item \texttt{tfsec-comment.yml} — Posts security findings as PR comments
\end{itemize}

Additionally, the DigitalOcean directory includes existing \texttt{README.md} with doctl and pass integration instructions for manual operations.

\section{Testing}

Terraform configurations are validated through:
\begin{itemize}
    \item Go-based integration tests (\texttt{terraform/test/main\_test.go})
    \item Plan dry-run validation via GitHub Actions (markdown.yml workflow)
    \item Manual state reconciliation checks (compare terraform state against live infrastructure)
\end{itemize}

State drift detection is recommended as a regular maintenance task — compare \texttt{terraform state list} output against actual cloud resource inventories quarterly.
 in lab-manual.tex after the Ansible chapter

\chapter{Terraform Infrastructure Management}

Terraform configurations in this lab manage cloud infrastructure across two providers: DigitalOcean and Google Cloud Platform. All configurations are version-controlled and include automated security scanning via GitHub Actions.

\section{DigitalOcean Infrastructure}

The \texttt{terraform/digital\_ocean/} directory manages the external-facing web presence for bitsmasher.net.

\subsection{Provisioned Resources}

\begin{itemize}
    \item \textbf{Droplet:} 512MB Debian 12 in lon1 region (www.bitsmasher.net)
    \item \textbf{Domain:} bitsmasher.net with full DNS management
    \item \textbf{DNS Records:} A, MX, TXT (SPF/DMARC/DKIM), NS, and ACME challenge records
\end{itemize}

\subsection{Key DNS Records}

\begin{center}
\begin{tabular}{lll}
\hline
Name & Type & Purpose \\ \hline
www.bitsmasher.net & A & Web server host (178.62.60.55) \\
@ & MX & mail.protonmail.ch 10 (ProtonMail routing) \\
@ & TXT & SPF and DMARC policies \\
protonmail.\_domainkey & CNAME & DKIM email verification \\
\_acme-challenge.*.lab & TXT & Let's Encrypt certificate challenges \\ \hline
\end{tabular}
\end{center}

\subsection{Configuration Notes}

The droplet includes:
\begin{itemize}
    \item Automated backups enabled (immutable against accidental deletion)
    \item SSH key injection via terraform data source
    \item No external monitoring (consider enabling DO agent for uptime tracking)
\end{itemize}

Provider authentication uses a DigitalOcean API token managed through pass/direnv integration. The token is passed at runtime, never committed to the repository.

\section{Google Cloud Platform Infrastructure}

The \texttt{terraform/google/} directory manages GCP resources for Stash House free-tier operations.

\subsection{Resources Managed}

\begin{itemize}
    \item \textbf{VPC Network:} \texttt{stash-house-vpc} dedicated network
    \item \textbf{Compute Instance:} e2-micro Debian 12 (Always Free Tier eligible) in us-central1-a
    \item \textbf{Firewall Rule:} \texttt{allow-ssh} opening port 22 to the VPC
    \item \textbf{Billing Budget Alert:} Tripwire alert at \$0.01 threshold for free-tier monitoring
\end{itemize}

The GCP configuration is intentionally minimal, designed around Google's Always Free Tier limits to maintain zero-cost infrastructure.

\section{Security Scanning}

All Terraform configurations are automatically scanned by two GitHub Actions workflows:
\begin{itemize}
    \item \texttt{tfsec-sarif.yml} — IaC security scanning (generates SARIF reports)
    \item \texttt{tfsec-comment.yml} — Posts security findings as PR comments
\end{itemize}

Additionally, the DigitalOcean directory includes existing \texttt{README.md} with doctl and pass integration instructions for manual operations.

\section{Testing}

Terraform configurations are validated through:
\begin{itemize}
    \item Go-based integration tests (\texttt{terraform/test/main\_test.go})
    \item Plan dry-run validation via GitHub Actions (markdown.yml workflow)
    \item Manual state reconciliation checks (compare terraform state against live infrastructure)
\end{itemize}

State drift detection is recommended as a regular maintenance task — compare \texttt{terraform state list} output against actual cloud resource inventories quarterly.


\end{document}
